% 日本語 LuaLaTeX。2026-11-22 / v0.2
% エンジン、出力先、SyncTeX は repo ルートの .latexmkrc に従う。
\RequirePackage{fix-cm}
\documentclass[a4paper,11pt]{ltjsarticle}
\usepackage{amsmath,amssymb,amsthm,mathtools,mathrsfs}
\DeclareFontFamily{U}{rsfs}{\skewchar\font127}
\DeclareFontShape{U}{rsfs}{m}{n}{<-6> rsfs5 <6-8> rsfs7 <8-> rsfs10}{}
\usepackage[margin=25mm]{geometry}
\usepackage{booktabs,tabularx,array,enumitem}
\usepackage[unicode,pdfusetitle,colorlinks=true,
  linkcolor=blue,citecolor=blue,urlcolor=blue]{hyperref}
\numberwithin{equation}{section}
\theoremstyle{plain}
\newtheorem{theorem}{定理}[section]
\newtheorem{proposition}[theorem]{命題}
\newtheorem{lemma}[theorem]{補題}
\newtheorem{corollary}[theorem]{系}
\theoremstyle{definition}
\newtheorem{definition}[theorem]{定義}
\newtheorem{assumption}[theorem]{仮定}
\newtheorem{example}[theorem]{例}
\newtheorem{problem}[theorem]{問題}
\theoremstyle{remark}
\newtheorem{remark}[theorem]{注意}
\renewcommand{\proofname}{証明}
\newcommand{\kk}{k}
\newcommand{\OO}{\mathcal O}
\newcommand{\PP}{\mathbb P}
\newcommand{\Gm}{\mathbb G_m}
\newcommand{\HH}{\mathcal H_2}
\newcommand{\HHstack}{\mathfrak H_2}
\DeclareMathOperator{\Hom}{Hom}
\DeclareMathOperator{\End}{End}
\DeclareMathOperator{\Ext}{Ext}
\DeclareMathOperator{\Aut}{Aut}
\DeclareMathOperator{\GL}{GL}
\DeclareMathOperator{\PGL}{PGL}
\DeclareMathOperator{\Gr}{Gr}
\DeclareMathOperator{\Spec}{Spec}
\DeclareMathOperator{\rank}{rank}
\DeclareMathOperator{\Pic}{Pic}
\DeclareMathOperator{\Jac}{Jac}
\DeclareMathOperator{\coh}{coh}
\DeclareMathOperator{\qgr}{qgr}
\DeclareMathOperator{\Perf}{Perf}
\DeclareMathOperator{\Sym}{Sym}
\DeclareMathOperator{\coker}{coker}
\DeclareMathOperator{\im}{im}
\DeclareMathOperator{\Supp}{Supp}
\title{非可換二次曲面の八本・十本・十二本矢表示：\\
  モジュライの開部分、反標準 avatar、内在的忘却}
\author{}
\date{2026年11月22日\quad v0.2（公表用草稿）}
\begin{document}
\maketitle
\begin{abstract}
非可換二次曲面の四頂点表示について、矢数の変化をモジュライの族として記述する。
有限 foundation の生成矢、関係の双対を加えた箙、非最小な potential 表示を区別する。
固定した Hom 次元を持つ有限有向代数の結合律スキームを構成し、
二つの隣接乗法の階数から、八本、二種類の十本、十二本への層別を得る。
少ない矢数の開部分は、正確に十本または十二本を要求する層ではなく、
非最小表示を含めた共通モジュライの開部分である。
重み四 potential の二つの二次係数による消去は任意の基底上で明示できるが、
通常の right-equivalence 商には消去される対の自己同型が残る。
従って Jacobian 代数の一致を、そのまま表示スタックの同値とすることはできない。

幾何側では、$\Sigma_0$ または $\Sigma_2$ 上の反標準因子と
非自明な $\Pic^0$ パラメータを用い、直積型 square avatar と
二種類の十本矢 chart の関係を記述する。
統一偏極と rotated square の Hilbert 偏極は異なり、
$\Sigma_2$ の出現だけでは最小十二本矢型を意味しない。
最後に、構造対象のみを指定した強い内在的モジュライ $\HHstack$ への忘却を扱う。
相対 heart、反 Serre の豊富な exact 平方根、同じ根の全軌道環の AS4 被覆を
確認した部分では忘却射が定義される。
この適合部分の開性や、統一 avatar と全相対 AS 代数とのスタック同値は
別の認識問題として残し、有限表示の代数性とは区別する。
\end{abstract}
\tableofcontents

\section{Introduction}
\label{sec:qa-introduction}

三次元 AS 正則 $\mathbb Z$-代数の四周期表示では、同じ幾何的構成に対して
八本、十本、十二本の矢を持つ箙が現れる。
これらを別々の分類集合として並べるだけでは、
生成元と関係が同時に消去される変形や、異なる foundation の比較を捉えられない。
本稿は、まず共通の有限代数のモジュライを構成し、その上の二つの開条件と
二つの階数低下条件として矢数を整理する。
次に反標準 avatar の幾何的 chart を比較し、
最後に構造対象付き Abel 圏への忘却に必要な条件を分離する。

\subsection{三つの問題と本稿の結果}

本稿で扱う問題には、次の三つの水準がある。
\begin{enumerate}
\item \emph{有限表示のモジュライ。}
  四頂点間の Hom 次元と結合律を固定した代数スタックは有限型 Artin である。
  隣接する二つの乗法が両方とも同型である部分、一方が同型である二つの部分、
  双方に一次元の欠損を許す部分を比較する。
  これが八本・十本・十二本の開部分の正確な有限代数版である。
\item \emph{AS 代数と avatar。}
  四周期の全 AS 条件と幾何的 avatar の対応には、有限 foundation の分類に加えて
  全単純分解、noetherian 性、相対平坦性、降下が必要である。
  体上の既存原稿の分類結果と、今回の有限代数・幾何的 chart の証明を区別する。
\item \emph{内在的な対象。}
  構造対象 $\OO$ だけを指定した $\HHstack$ は、選んだ箙、potential、
  foundation、平方根、ambient を忘れる。
  しかしそれらの存在条件を失うわけではない。
  特に AS3 正則性や $\Sigma_2$ という有限表現モジュライだけから、
  $\HHstack$ の全相対条件は導けない。
\end{enumerate}

有限代数版では、二種類の十本矢 chart の交わりが八本矢 chart である。
二つの十本矢 chart の和集合の補集合が最小十二本矢の層となる。
これに対し、各 chart の内部で正確に十本が最小生成矢として残る部分は
八本矢部分の境界である。
このため、本稿では上限を許す chart と最小本数が一定の層を同じ記号で表さない。

直積型 collection
\[
 (\OO,\OO(1,0),\OO(0,1),\OO(1,1))
\]
そのものの有限箙は八本、四本の最小関係の双対を加えた relation QP は十二本である。
固定した一頂点 mutation により、八本と一方の十本を含む chart が得られる。
従って直積型と十本矢 chart の対応は mutation を介した対応であり、
元の箙の矢数を十本と数えるという意味ではない。

\subsection{先行研究と依存関係}

有限箙の関係式を moduli stack として扱う一般的枠組みは
Abdelgadir--Okawa--Ueda \cite[Definitions 3.5, 4.1; Propositions 3.6, 4.2]{AOU2014}
による。本稿は四頂点の Hom 次元を固定して結合律を具体的な行列方程式で書き、
二つの乗法の階数が最小箙に与える影響を計算する。
これは一般論の新しい証明を主張するものではない。
Derksen--Weyman--Zelevinsky \cite[Theorem 4.6; Section 5]{DWZ2008}
は完備道代数の potential の splitting と mutation を構成した。
ここで必要な消去は内部次数四のため有限回の三角変換で済み、
通常の道代数にも任意基底にも同じ式を適用できる。
ただし pair の消去と quotient stack の慣性の消去は別の操作である。

AS 正則 quadric の cubic 型については Van den Bergh \cite{VdB2011}、
その幾何的データと potential のモジュライについては
Okawa--Ueda \cite{OU2014} が基礎となる。
ここで用いる square・rotated square の反標準 avatar と十本矢 mutation は
大川--植田の日本語草稿 v7 \cite{AvatarVSeven2026} に、
十二本矢型の追加と統一 avatar は研究ノート \cite{Q12Avatar2026} に記載されている。
幾何的な相対 twist は \cite{Q12Weyl2026}、
これらの関係の整理は \cite{Synthesis2026} を参照する。
これらは未公刊の研究資料であり、そこに記載された全 AS 分類の証明を
今回独立に検証済みの結果として引用しない。
本稿では必要な有限線形代数・消去・chart の証明を本文に置き、
追加の再構成入力を要する主張を明記する。

内在的な $\HH$ とその相対化 $\HHstack$ は、同じ研究の主論文
\cite{QuadricIntrinsic2026} の対象である。
分類主論文は構造対象を保つ exact 同値による分類を扱い、
新しい相対モジュライ原稿 \texttt{nc-quadric-moduli.tex} は
相対認識と境界の幾何を担当する（現在は証明移管中）。
本稿は表示の層別と avatar の幾何を独立に説明する。
主論文中の全分類や相対再構成の証明は本稿に重複させない。
本稿の結果も、全 $\HHstack$ の代数性、全 avatar の相対 AS 認識、
または公刊に向けた包括的な新規性確認の完了を意味しない。

\subsection{規約}

$\kk$ は標数零の代数閉体とする。
族は $\kk$ 上局所有限型スキームの fppf site 上で考え、非被約基底を含む。
矢の積は合成順で、$ab$ は先に $b$、次に $a$ を通る道である。
有限 foundation は四つの順序付き頂点を保持し、
右加群規約に対応する導来同値と標準 heart の exact 同値を区別する。
基底上の係数を消去に使う条件は可逆性であり、単に零でないことではない。

% !TEX root = ../quiver-avatar-moduli.tex
% 2026-11-21 / v1.0
\section{矢の本数、安定化表示、および開部分スタック}
\label{sec:qa-presentations}

本節では、有限 foundation の最小矢、relation-dual の戻り矢を加えた
箙、重み付き箙の非最小表示を区別する。
この区別のもとで、八本矢のモジュライが十本矢のモジュライの開部分となる
という主張を、任意の底上で正確に定式化する。
基礎体を $\kk$ とし、積 $ba$ は $a$ の後に $b$ を合成する規約とする。
本節の代数的な消去には標数の制限を要しない。

\subsection{直積型 collection の最小表示}

\begin{proposition}[直積型 collection の矢と関係]
\label{prop:qa-product-count}
$X=\PP^1\times\PP^1$ とし、
\[
 (E_0,E_1,E_2,E_3)
   =(\OO,\OO(1,0),\OO(0,1),\OO(1,1))
\]
と置く。有限代数 $B=\End_X(\bigoplus_{i=0}^3E_i)$ の最小箙は
\[
 a_i:0\longrightarrow1,\quad b_j:0\longrightarrow2,\quad
 u_j:1\longrightarrow3,\quad v_i:2\longrightarrow3
 \qquad(i,j=0,1)
\]
の八本の矢を持ち、最小関係は
\begin{equation}
\label{eq:qa-square-relations}
 r_{ij}=u_ja_i-v_ib_j\qquad(i,j=0,1)
\end{equation}
の四本である。従って relation-dual の戻り矢
$c_{ij}:3\to0$ を加えた relation QP は十二本矢であり、potential は
\begin{equation}
\label{eq:qa-square-potential}
 W_{\mathrm{sq}}=\sum_{i,j=0}^1c_{ij}(u_ja_i-v_ib_j)
\end{equation}
である。これは道長に関して reduced である。
\end{proposition}

\begin{proof}
$V=H^0(X,\OO(1,0))$、$U=H^0(X,\OO(0,1))$ はともに二次元である。
非恒等 Hom のうち、$\Hom(E_0,E_3)$ 以外の非零成分は
$V,U,U,V$ の四つである。
$\Hom(E_1,E_2)=H^0(X,\OO(-1,1))=0$ であり、
$\Hom(E_0,E_3)=V\otimes U$ への二つの道の合成はともに同型である。
従って二段階の道空間 $(U\otimes V)\oplus(V\otimes U)$ から
$V\otimes U$ への核は四次元で、\eqref{eq:qa-square-relations} が基底となる。
それより長い道は存在せず、これらが関係の全てである。
関係に一次部分はないので八本の矢は最小であり、
関係を他の矢で左右から乗じて新たな関係を得る余地もない。
四本の双対矢を加えると十二本となる。
\eqref{eq:qa-square-potential} は道長三なので、その二次部分は零である。
\end{proof}

\begin{remark}[十二本という数だけでは箙を指定できない]
\label{rem:qa-two-twelve-quivers}
命題~\ref{prop:qa-product-count} の箙を $Q_{\mathrm{sq}}$ と書く。
後述する $Q_{12}$ は四辺に各二本の矢を持つ有向四角形と四本の対角矢からなる。
$Q_{\mathrm{sq}}$ は $3\to0$ に四本の矢を持ち、$1\to2$ の矢を持たない。
従って両者は同じ十二本矢の箙ではない。
また、直積型 collection 自身の relation QP を十本矢と数えることはできない。
十本矢との対応を主張する際には、mutation と collection の変更を明示する必要がある。
この数え方は \cite[Remark 38.2]{AvatarVSeven2026} と一致する。
\end{remark}

\subsection{二つの可縮 pair を持つ共通表示}

四頂点を $\mathbb Z/4\mathbb Z$ で標識し、$Q_8$ を
各 $i\to i+1$ に二本の矢を持つ有向四角形とする。
その矢の重みを一とする。
$Q_{10}^{\mathrm{odd}}$ はこれに
\[
 x:1\longrightarrow3,\qquad y:3\longrightarrow1
\]
を加えた箙で、追加矢の重みは二とする。
$Q_{10}^{\mathrm{even}}$ は代わりに
\[
 u:0\longrightarrow2,\qquad v:2\longrightarrow0
\]
を加えたもの、$Q_{12}$ は両方の pair を加えたものとする。
以下の potential は重み四の巡回同値類であり、potential の定数倍は
同一視しない。

\begin{lemma}[共通 potential の形]
\label{lem:qa-universal-potential}
$Q_{12}$ の重み四の potential は一意に
\begin{equation}
\label{eq:qa-universal-potential}
 W=H+yq+px+vq'+p'u+\lambda yx+\mu vu
\end{equation}
と書ける。ここで $H$ は $Q_8$ の道長四の potential、
$q:1\to3$、$p:3\to1$、$q':0\to2$、$p':2\to0$ は
$Q_8$ の道長二の線形結合であり、$\lambda,\mu$ は係数である。
従って基底を固定した potential 空間は
\[
 P_8\simeq\mathbb A^{16},\qquad
 P_{10}^{\mathrm{odd}}\simeq P_{10}^{\mathrm{even}}\simeq\mathbb A^{25},
 \qquad P_{12}\simeq\mathbb A^{34}
\]
である。これらの記述は任意の $\kk$-scheme 上で成立する。
\end{lemma}

\begin{proof}
重み四の cycle が重み二の矢を持たなければ、四辺を一周するので
$2^4=16$ 個の道が巡回同値類の基底を与える。
重み二の矢を一つ持つ場合は、残りの二本が $Q_8$ の道長二の道となり、
四本の対角矢それぞれについて四個の可能性がある。
重み二の矢を二つ持つ cycle は $yx$ と $vu$ の二つだけである。
この三種類は巡回回転で混ざらないので、直和分解と次元の式を得る。
同じ道の基底を任意の可換係数環上で用いれば相対的な主張も従う。
\end{proof}

\begin{proposition}[任意の底上の有限な消去]
\label{prop:qa-relative-elimination}
$T$ 上で \eqref{eq:qa-universal-potential} が与えられているとする。
$\lambda\in\Gamma(T,\OO_T)^\times$ なら、
\[
 x'=x+\lambda^{-1}q,\qquad y'=y+\lambda^{-1}p
\]
によって
\begin{equation}
\label{eq:qa-one-pair-elimination}
 W=\lambda y'x'+H-\lambda^{-1}pq+vq'+p'u+\mu vu
\end{equation}
となる。従って通常の重み付き道代数で定義した Jacobian 代数について
\[
 \Jac(Q_{12},W)
 \simeq
 \Jac(Q_{10}^{\mathrm{even}},H-\lambda^{-1}pq+vq'+p'u+\mu vu)
\]
が成り立つ。$\mu$ が可逆な場合も同様で、両方が可逆なら
\begin{equation}
\label{eq:qa-two-pair-elimination}
 \Jac(Q_{12},W)
 \simeq\Jac(Q_8,H-\lambda^{-1}pq-\mu^{-1}p'q').
\end{equation}
これらの構成は任意の基底変換と両立する。
\end{proposition}

\begin{proof}
$p,q$ は $x,y,u,v$ を含まないので、表示した変換は逆も多項式である
重み付き道代数の自己同型である。式は積を展開して確認できる。
変換後の cyclic derivative は $\lambda x'$ と $\lambda y'$ を含むため、
Jacobian 商では $x'=y'=0$ である。
残りの cyclic derivative は残余 potential のものに一致する。
二つの pair の変換は互いの矢を変更せず、順序を問わず実行できる。
用いた演算は係数の加減乗除と、表示した可逆元の逆元だけである。
従って同型は任意の基底変換後にも同じ式で成立する。
\end{proof}

\begin{remark}[一般の splitting theorem との関係]
一般の QP の splitting theorem は完備道代数上の
\cite[Theorem 4.6]{DWZ2008} による。
命題~\ref{prop:qa-relative-elimination} では混合項が追加矢を含まないため、
有限回の三角変換だけで分離できる。従って完備化も極限操作も不要である。
この相対計算は AS 正則性や各次数成分の平坦性を主張していない。
それらは別の仮定または定理を必要とする。
\end{remark}

幾何点上の最小矢数は次の四つの場合に分かれる。
表の $\lambda=0$、$\mu=0$ は、その幾何点の剰余体での等式である。
\begin{center}
\begin{tabular}{c|c|c}
係数 & 残る最小対角矢 & 最小矢数 \\
\hline
$\lambda\ne0,\ \mu\ne0$ & なし & $8$ \\
$\lambda=0,\ \mu\ne0$ & $x,y$ & $10$（odd）\\
$\lambda\ne0,\ \mu=0$ & $u,v$ & $10$（even）\\
$\lambda=0,\ \mu=0$ & $x,y,u,v$ & $12$
\end{tabular}
\end{center}
消去後の potential は道長三以上なので、Jacobian ideal に一次関係はない。
これが表の最小性の根拠である。
非被約な底上では、全幾何点で係数が零でも係数自身が零とは限らない。
例えば $\lambda=\epsilon$ on $\Spec\kk[\epsilon]/(\epsilon^2)$ は可逆でなく、
最後の pair を消去できない。

\subsection{通常の QP スタックに残る慣性}

ここでは頂点を固定する全ての重み付き道代数自己同型を射とする。
$G_8=(\GL_2)^4$ と置くと、重み一の矢の像は基底変更で尽くされる。
十本矢では、さらに各追加矢の非零スカラー倍と、
同じ始点・終点を持つ道長二の道の付加が許される。従って
\[
 G_{10}=(G_8\times\Gm^2)\ltimes\mathbb G_a^8,
 \qquad
 G_{12}=(G_8\times\Gm^4)\ltimes\mathbb G_a^{16}.
\]
群の作用は potential の変数変換であり、
$\mathcal P_r=[P_r/G_r]$ を通常の重み付き QP 商スタックとする。
この節では頂点スカラーの rigidification は行わない。

\begin{proposition}[可縮 pair の残余慣性]
\label{prop:qa-stabilized-inertia}
上の群と頂点標識を用いると、
\begin{align}
 [\{\lambda\ne0\}\subset P_{10}^{\mathrm{odd}}/G_{10}]
   &\simeq \mathcal P_8\times B\Gm,\label{eq:qa-qp-gerbe-one}\\
 [\{\lambda\mu\ne0\}\subset P_{12}/G_{12}]
   &\simeq \mathcal P_8\times B(\Gm^2).
   \label{eq:qa-qp-gerbe-two}
\end{align}
特に、消去によってこれらの通常の QP スタックを
$\mathcal P_8$ と同定することはできない。
\end{proposition}

\begin{proof}
十本矢の場合、命題~\ref{prop:qa-relative-elimination} により
$p=q=0$ とでき、さらに一方の矢のスケールを変えて $\lambda=1$ とできる。
これは係数環上の可逆元の逆を用いる代数的な正規化である。
正規化した切片は $H+yx$ の形の potential の空間 $P_8$ である。

この切片の二点 $H+yx$ と $\widetilde H+yx$ を結ぶ自己同型を調べる。
重み一部分の作用を $g\in G_8$、対角矢の作用を
$x\mapsto\alpha x+q_0$、$y\mapsto\beta y+p_0$ と書く。
変換後に $yq_0$ と $p_0x$ の混合項が消えるには、
$\alpha,\beta$ が可逆なので $p_0=q_0=0$ が必要である。
$yx$ の係数から $\alpha\beta=1$ を得る。
従って切片間の全ての射は
\[
 (g,t):\quad H\longmapsto gH,\quad
 x\longmapsto tx,\quad y\longmapsto t^{-1}y
\]
であり、その群は $G_8\times\Gm$ である。
後半の $\Gm$ は $P_8$ に自明に作用する。
各対象が fppf 局所にこの切片へ移り、切片間の射がちょうど上の群で
記述されるので、商スタックの同値~\eqref{eq:qa-qp-gerbe-one} を得る。
二つの pair について同じ正規化を独立に行えば
\eqref{eq:qa-qp-gerbe-two} を得る。
ここでは切片とその安定化群を明示したので、単に gerbe と述べるだけでなく
直積との同値まで得られている。
\end{proof}

\begin{remark}[表示を忘れる操作]
消去される pair の変換 $x\mapsto tx$、$y\mapsto t^{-1}y$ は
Jacobian 商では恒等である。
従って通常の QP groupoid から代数 groupoid への忘却は忠実でない。
同じ障害は \cite[Proposition 52.3]{AvatarVSeven2026} で指摘されている。
一般に関係式の表示の自己同型と商代数の自己同型を区別すべきことは
\cite[Proposition 3.3]{AOU2014} にも現れる。
有限有向箙の関係式スタックの有限型性
\cite[Definition 3.5, Proposition 3.6]{AOU2014} は有用であるが、
その定理だけから、ここで扱う無限次元 graded 代数の全次数平坦な族や
内在的モジュライの代数性は従わない。
\end{remark}

\subsection{代数の射を用いた正確な開埋込み}

\begin{definition}[代数の fppf スタック]
\label{def:qa-algebra-stack}
$\mathfrak M_{12}^{\mathrm{alg}}$ の $T$ 上の対象は、
$A_0=\OO_T^4$ と四つの標識付き冪等元を持つ非負次数付き代数
$A=\bigoplus_{d\ge0}A_d$ で、各 $A_d$ が有限階数局所自由、かつ
fppf 局所に $\Jac(Q_{12},W)$ の形の表示を持つものとする。
射は頂点と次数を保つ全ての代数同型とし、potential、矢の基底、
表示からの全射はデータに含めない。
$\mathfrak M_{10}^{\mathrm{odd,alg}}$、
$\mathfrak M_{10}^{\mathrm{even,alg}}$、$\mathfrak M_8^{\mathrm{alg}}$
もそれぞれの箙を用いて同様に定義する。
これらは fppf 降下を持つ代数のスタックの充満部分スタックである。
この定義に AS 正則性は含めない。
\end{definition}

各対象について、添字を $\mathbb Z/4\mathbb Z$ で読み、有限表示層
\begin{equation}
\label{eq:qa-indecomposables}
 C_i(A)=\operatorname{coker}\left(
 e_{i+2}A_1e_{i+1}\otimes_{\OO_T}e_{i+1}A_1e_i
 \longrightarrow e_{i+2}A_2e_i\right)
\end{equation}
を定める。これは重み二における新しい生成元の層であり、
定義は表示に依存しない。

\begin{theorem}[矢の少ない代数スタックの開性]
\label{thm:qa-open-algebra-stacks}
可縮 pair を追加することによって、充満な開埋込み
\begin{equation}
\label{eq:qa-open-diagram}
 \begin{aligned}
 \mathfrak M_8^{\mathrm{alg}}&\lhook\joinrel\longrightarrow
 \mathfrak M_{10}^{\mathrm{odd,alg}}
 \lhook\joinrel\longrightarrow\mathfrak M_{12}^{\mathrm{alg}},\\
 \mathfrak M_8^{\mathrm{alg}}&\lhook\joinrel\longrightarrow
 \mathfrak M_{10}^{\mathrm{even,alg}}
 \lhook\joinrel\longrightarrow\mathfrak M_{12}^{\mathrm{alg}}.
 \end{aligned}
\end{equation}
を得る。
両十本矢開部分の交わりは $\mathfrak M_8^{\mathrm{alg}}$ である。
正確には、それぞれの像は
\begin{align*}
 \mathfrak M_{10}^{\mathrm{odd,alg}}&=\{C_0=C_2=0\},\\
 \mathfrak M_{10}^{\mathrm{even,alg}}&=\{C_1=C_3=0\},\\
 \mathfrak M_8^{\mathrm{alg}}&=\{C_0=C_1=C_2=C_3=0\}
\end{align*}
である。開埋込みという主張は任意の scheme からの基底変換について成立し、
各スタックの代数性を仮定しない。
\end{theorem}

\begin{proof}
十本矢の表示に $vu$ を追加しても Jacobian 代数は変わらず、
八本矢からは $yx$ と $vu$ を追加できる。
従って同じ代数を対象として送る関手が存在する。
射は両辺とも全ての次数付き頂点固定代数同型なので、関手は充満忠実である。

局所表示~\eqref{eq:qa-universal-potential} で
\eqref{eq:qa-indecomposables} を計算する。
重み二の関係は四本の重み二の矢による cyclic derivative だけである。
道長二の積を割ると、これらの関係は
$\lambda x=\lambda y=\mu u=\mu v=0$ となる。従って局所的に
\[
 C_1\simeq C_3\simeq\OO_T/(\lambda),\qquad
 C_0\simeq C_2\simeq\OO_T/(\mu).
\]
特に $C_0=C_2=0$ と $\mu$ の可逆性は同値であり、
命題~\ref{prop:qa-relative-elimination} により
$Q_{10}^{\mathrm{odd}}$ の表示を得る。他の二つの場合も同様である。
逆に小さい箙の表示では、対応する生成元層は明らかに零である。
これにより像の表示を得る。

有限表示層 $C_i$ の消滅 locus は、その support の補集合という開集合である。
また cokernel は任意の基底変換と両立する。
従って任意の $T\to\mathfrak M_{12}^{\mathrm{alg}}$ による各関手の
基底変換は、対応する $C_i$ の消滅で定まる $T$ の開部分そのものである。
これは開埋込みの主張と交わりの式を同時に証明する。
\end{proof}

\begin{remark}[最小型、相対型、および追加条件]
\label{rem:qa-open-limits}
十本矢の開部分に八本矢の対象が入るのは、
定義~\ref{def:qa-algebra-stack} が非最小の十本矢表示も許すためである。
幾何点で最小矢がちょうど十本の対象だけを取れば、八本矢の対象は含まれない。
相対的に固定した最小型を扱うときは、残る $C_i$ が階数一の局所自由層で
あることを課す。この条件は局所表示では対応する mass 係数が底上で
零となることを要求し、全幾何点で零という条件より強い。

本節の開性は、同型と基底変換に不変な追加条件を課した共通の
部分スタックへ制限しても保たれる。
ただし AS 条件を満たす部分がどの部分であるか、その部分の代数性、
avatar からの全域被覆、さらに構造対象を保つ内在的圏への射は
この定理の結論に含まれない。
また $\mathfrak M_r^{\mathrm{alg}}$ では四頂点と周期的な graded 表示を
保持している。これを忘れた $\mathcal O$-marked の exact 同値を、
次数付き代数同型と同一視してはならない。
\end{remark}

% !TEX root = ../quiver-avatar-moduli.tex
% 2026-11-21
\section{有限 foundation の二つの十本矢開部分}
\label{sec:qa-finite-foundations}

本節では、矢の本数による開部分を実際の代数スタックとして構成する。
対象は四頂点の有限 directed algebra であり、AS 正則性や標準 heart は
この構成の条件に含めない。最後に数える矢は、有限代数そのものの矢と
最小関係に双対な戻り矢を合わせた relation completion の矢である。
従って、ここで得る開部分の定理と、AS 代数・avatar の相対分類定理とは
異なる主張である。

\subsection{全ての積を保持する有限表示}

\(U,V,W\) を二次元、\(E,F\) を四次元、\(N\) を六次元の
\(\kk\)-ベクトル空間とし、頂点成分を
\[
 A_{01}=U,\quad A_{12}=V,\quad A_{23}=W,\qquad
 A_{02}=E,\quad A_{13}=F,\quad A_{03}=N
\]
と固定する。対角成分は \(A_{ii}=\kk\)、逆向き成分は零とする。
非自明な積は四写像
\[
 \begin{aligned}
 a&:V\otimes U\longrightarrow E,&
 b&:W\otimes V\longrightarrow F,\\
 c&:W\otimes E\longrightarrow N,&
 d&:F\otimes U\longrightarrow N
 \end{aligned}
\]
であり、結合律はただ一つの写像等式
\begin{equation}
 c(1_W\otimes a)=d(b\otimes1_U)
 \quad\text{in }\Hom(W\otimes V\otimes U,N)
 \label{eq:qa-associativity}
\end{equation}
である。この等式の定める閉アフィンスキームを \(Z\) と書く。
\(Z\subset\mathbb A^{128}\) は有限型である。
全ての積を保持するため、後に \(a\) または \(b\) の階数が落ちても
有限代数のベクトル空間としての階数二十四は変わらない。

\begin{definition}[foundation の許容開部分]
\label{def:qa-foundation-open}
\(Z\) の開部分 \(Z^{\mathrm{nd}}\) を次の条件で定める。
\begin{enumerate}
\item \((c,d): (W\otimes E)\oplus(F\otimes U)\to N\) は全射である。
\item 全ての幾何ファイバーで、\(\ker a\) と \(\ker b\) は
非零の分解可能テンソルを含まない。
\end{enumerate}
ここで \(Z^{\mathrm{nd}}\) の記号は上記の有限条件だけを意味し、
AS 正則性を意味しない。
\end{definition}

\begin{lemma}
\label{lem:qa-foundation-open-ranks}
定義~\ref{def:qa-foundation-open} は基底変更と両立する開条件である。
各幾何点で \(\rank a,\rank b\in\{3,4\}\) である。
\end{lemma}
\begin{proof}
全射性は最大 minor の非消滅条件である。
\(a\) について、\(Z\times\PP(V)\times\PP(U)\) 上で
\(a(v\otimes u)=0\) となる閉 incidence を取る。
射影空間の積からの射影は proper なので、その \(Z\) への像は閉である。
その補集合が第二条件を表す。\(b\) についても同じである。

二次元空間のテンソル積の射影化 \(\PP^3\) において、
分解可能テンソルは Segre 二次曲面をなす。
代数閉体上の射影直線はこの二次曲面と交わる。
従って二次元以上の核は非零の分解可能テンソルを含むので、
許される核の次元は零または一である。
\end{proof}

\[
 G=\GL(U)\times\GL(V)\times\GL(W)
       \times\GL(E)\times\GL(F)\times\GL(N)
\]
を各空間の基底変更で作用させる。
頂点の対角単元 \((\lambda_0,\lambda_1,\lambda_2,\lambda_3)\)
が \(A_{ij}\) 上で \(\lambda_j\lambda_i^{-1}\) 倍する作用の像を
\(K\subset G\) と書く。
共通スカラーを除けば \(K\simeq\Gm^3\) であり、
\(K\) は \(Z\) 上自明に作用する中心部分群である。

\begin{proposition}[有限代数の代数スタック]
\label{prop:qa-foundation-stack}
\[
 \mathfrak F^{\mathrm{lin}}_{\le12}=[Z^{\mathrm{nd}}/G],
 \qquad
 \mathfrak F_{\le12}=[Z^{\mathrm{nd}}/(G/K)]
\]
は有限型 Artin スタックである。
前者は上記の階数の頂点成分と積を持つ有限代数のスタックであり、
後者は頂点ごとの比較同型のスカラーによる中心 \(K\) を剛体化したものである。
この命題ではスタックの滑らかさも AS locus の開性も主張しない。
\end{proposition}
\begin{proof}
各頂点成分を局所自由層とする同じデータは、fppf 局所に基底を選ぶと
\(Z^{\mathrm{nd}}\) の点で表される。
代数の同型はこれらの空間の基底変更と正確に一致する。
従って前者の商は指定した moduli を表す。
\(G\) と \(G/K\) は滑らかな有限型代数群なので、
両商は有限型 Artin スタックである。
\(K\) は各対象に中心自己同型として含まれるため、
\([Z^{\mathrm{nd}}/G]\to[Z^{\mathrm{nd}}/(G/K)]\) はその剛体化である。
\end{proof}

\subsection{最小生成元と最小関係の数}

\begin{proposition}[八・十・十二本の有限代数的判定]
\label{prop:qa-foundation-arrow-count}
\(Z^{\mathrm{nd}}\) の幾何点に対応する代数を \(A\) とし、
\[
 \delta_a=4-\rank a,\qquad \delta_b=4-\rank b
 \quad\in\{0,1\}
\]
と置く。\(A\) の最小 directed quiver は
\(6+\delta_a+\delta_b\) 本の矢を持ち、最小関係は
\(2+\delta_a+\delta_b\) 本である。
従って、最小関係ごとに一本の戻り矢を加えた relation completion は
\begin{equation}
 8+2(\delta_a+\delta_b)
 \label{eq:qa-foundation-arrow-count}
\end{equation}
本の矢を持つ。
\end{proposition}
\begin{proof}
隣接頂点間には各二本、合計六本の矢がある。
\(E/\operatorname{im}a\) と \(F/\operatorname{im}b\) の基底を取り、
それぞれ \(0\to2\)、\(1\to3\) の追加生成矢とする。
\((c,d)\) の全射性により、\(0\to3\) の全ての元は
\(0\to2\to3\) または \(0\to1\to3\) を経由する積で生成されるので、
\(0\to3\) の追加生成矢はない。
従って生成矢の数は \(6+\delta_a+\delta_b\) である。

\(0\to2\)、\(1\to3\) の最小関係は、それぞれ
\(K_a=\ker a\)、\(K_b=\ker b\) であり、その次元は
\(\delta_a,\delta_b\) である。
これらの関係を課す前の \(0\to3\) の道の空間は
\[
 (W\otimes V\otimes U)
 \oplus\bigl(W\otimes(E/\operatorname{im}a)\bigr)
 \oplus\bigl((F/\operatorname{im}b)\otimes U\bigr)
\]
という、次元 \(8+2\delta_a+2\delta_b\) の空間である。
短い関係の生成する部分空間は最初の直和成分内の
\[
 (W\otimes K_a)+(K_b\otimes U)
\]
であり、この和は直和である。
実際、両核が一次元で共通の非零テンソルがあれば、
それは \(w\otimes\kappa_a=\kappa_b\otimes u\) と書ける。
\(w\) に対して非零となる線形形式で縮約すると、
\(\kappa_a\) が分解可能になり、定義に反する。
一方の核が零の場合は明らかである。

従って短い関係を課した後の \(0\to3\) 成分は次元八である。
結合律~\eqref{eq:qa-associativity} により \(N\) への写像が定まり、
これは \((c,d)\) の全射性によって全射である。
残る核は二次元であり、ちょうど二本の新たな最小関係を与える。
添字次数は三であるが、追加矢があるときは道長二と三の項を含み得る。
既に短い関係の生成する部分で割っているため、これらは最小関係であり、
四頂点の有向性から、それ以外の始点・終点に新たな関係はない。
これで関係数と \eqref{eq:qa-foundation-arrow-count} を得る。
\end{proof}

添字の差による重みを有限代数の矢 \(i\to j\) に与え、
同じ始点・終点を持つ関係の戻り矢には重み \(4-(j-i)\) を与える。
関係 potential の各項は重み四となる。
隣接矢六本と \(0\to3\) の二つの関係の戻り矢は、
四周期の各隣接頂点を結ぶ二本ずつの重み一の矢を与える。
\(\delta_a=1\) なら \(0\to2\)、\(2\to0\) の重み二の矢を一組、
\(\delta_b=1\) なら \(1\to3\)、\(3\to1\) の重み二の矢を一組加える。
従って重み付きの箙自体は、八本矢、二種類の十本矢、十二本矢の
それぞれの配置に正確に一致する。
この配置の一致は、potential の復元や AS 正則性、
QP の同型と有限代数の同型の比較まで含む主張ではない。

\subsection{二つの開部分とその共通部分}

\begin{theorem}[矢数を上限とする開部分]
\label{thm:qa-foundation-open-diagram}
\(Z^{\mathrm{nd}}\) 上で \(a\)、\(b\) が可逆となる開部分を
それぞれ \(Z_a,Z_b\) と書く。
\[
 \begin{aligned}
 \mathfrak F_{\le10}^{(a)}&=[Z_b/(G/K)],&
 \mathfrak F_{\le10}^{(b)}&=[Z_a/(G/K)],\\
 \mathfrak F_8&=[(Z_a\cap Z_b)/(G/K)]
 \end{aligned}
\]
と置く。上付き \((a)\)、\((b)\) は、階数低下を許す写像を示す。
これらは \(\mathfrak F_{\le12}\) の開部分スタックであり、
\begin{equation}
 \mathfrak F_8
 =\mathfrak F_{\le10}^{(a)}
       \cap\mathfrak F_{\le10}^{(b)}.
 \label{eq:qa-foundation-chart-intersection}
\end{equation}
両十本矢開部分の合併は、最小 relation completion の矢数が
高々十である locus である。
その補集合は、幾何点上では \(\rank a=\rank b=3\) の
最小十二本矢 locus である。
\end{theorem}
\begin{proof}
可逆性は行列式の非消滅であり、\(G/K\) の作用で不変である。
不変開部分からの商の射は representable な open immersion である。
交わりと合併は同じ基底上で計算できる。
矢数との対応は命題~\ref{prop:qa-foundation-arrow-count} による。
\end{proof}

\begin{remark}[最小本数の層と開部分の区別]
\label{rem:qa-exact-arrow-strata}
例えば \(\mathfrak F_{\le10}^{(a)}\) は八本矢の点も含む。
最小十本矢の点だけを取った層は
\(\mathfrak F_{\le10}^{(a)}\setminus\mathfrak F_8\) であり、
八本矢の moduli を開部分として含むのはこの層ではない。
非被約基底について「階数三の層」と呼ぶときは、単なる幾何点条件でなく、
対応する大きな minor の消滅と三次 minor の可逆性を課した
scheme-theoretic な階数 locus を用いる。
本節の開部分は可逆性で定義されており、この区別による曖昧さを持たない。
また、選んだ最小関係とその戻り矢から作る QP の自己同型群を、
有限代数の自己同型群と同一視したわけではない。
QP から有限代数への回復と射の比較には別の証明が必要である。
\end{remark}

\subsection{cubic表示の幾何条件に対する次数六以下の明示的検査}
\label{subsec:qa-explicit-geometric-fourfold-test}

既存のnormal foundation chartで最初の積が同型となる場合、
論文B~\cite[\texttt{prop:cubic-foundation-open}]{QuadricModuli2026}
は二次元空間$V_0,V_1,V_2,V_3$の
四重テンソル$w\in V_0\otimes\cdots\otimes V_3$の
geometric性を使ってcubic逆構成を行う。
以下はこのgeometric開を
\emph{完全に明示的な次数四・六の多項式条件}に書き直す。

\begin{lemma}[四重テンソルのgeometric性のadjugate判定]
\label{lem:qa-explicit-cubic-geometric-adjugate}
$k$を代数閉体とし、各$V_i$を二次元として、
$w\in V_0\otimes V_1\otimes V_2\otimes V_3$
を取る。$i\in\mathbb Z/4\mathbb Z$について、
隣接二因子の縮約行列を
\[
 M_i:(V_i\otimes V_{i+1})^\vee
       \longrightarrow V_{i+2}\otimes V_{i+3}
\]
とする。基底を選んで$M_i$を$4\times4$行列とみなす。
$A_i=\operatorname{adj}(M_i)$をその古典的adjugateとし、
その第$j$列を二因子の$2\times2$テンソルに
reshapeした行列を$C_{ij}$（$1\le j\le4$）とする。
\[
 d_i=\det M_i,\qquad
 h_{ij}=\det C_{ij}
 \tag{Geo-poly}\label{eq:qa-cubic-geometric-minors}
\]
と置く。$d_i$は$w$の成分の次数四、
$h_{ij}$は次数六の多項式である。

Van den Berghのgeometric条件
\[
 (\phi_i\otimes\phi_{i+1})\mathbin{\lrcorner}w\ne0
 \quad\text{for all }i,\quad
 0\ne\phi_i\in V_i^\vee,\ \
 0\ne\phi_{i+1}\in V_{i+1}^\vee
\]
と次の有限条件は同値である。
\[
 \boxed{\quad
  \text{各$i=0,1,2,3$について、}
  \quad d_i\ne0
  \ \text{または}\ 
  h_{i1},h_{i2},h_{i3},h_{i4}
  \text{の少なくとも一つが非零}.\quad}
 \tag{Geo-test}\label{eq:qa-cubic-geometric-finite-test}
\]
従ってgeometric四重テンソル開は、
\emph{四つの次数四式と十六の次数六式だけ}
を使って定義される。
各$M_i$の行・列の基底変更に対し
この非消滅条件は不変なので、
有限foundationの相対chartへ降下する。

これは前段の一般存在定理
~\ref{cor:qa-realizable-foundation-finite-open-criterion}
で未計算だった\emph{全実現可能開}の方程式ではなく、
$F$が可逆であるcubic部分における
geometric条件の具体的な判定式である。
\end{lemma}
\begin{proof}
固定した$i$について、$K_i=\ker M_i$を
四次元空間$(V_i\otimes V_{i+1})^\vee$の部分空間とする。
geometric条件は、$K_i$に非零の
\emph{分解可能}二因子テンソルがないことと同値。
これは
$\PP(K_i)\cap\operatorname{Seg}(\PP^1\times\PP^1)=\varnothing$
という条件である。
Segre像は$\PP^3$の滑らかな二次曲面で、
分解可能な$2\times2$行列は行列式零によって特徴付けられる。

$\operatorname{rank}M_i=4$なら$K_i=0$なので条件は成立し、
$d_i\ne0$と同値である。
$\operatorname{rank}M_i\le2$なら
$\dim K_i\ge2$なので$\PP(K_i)$は少なくとも直線を含む。
代数閉体上で任意の射影直線は二次曲面と交わる
（線上で行列式を制限した二次形式に零点がある、
恒等零の場合も交わる）。
従ってgeometric条件は失敗する。
このとき全ての$3\times3$小行列式が零なので
$A_i=0$となり、全ての$h_{ij}$は零である。

残る$\operatorname{rank}M_i=3$の場合、
$A_i$は階数一の非零行列で、像が
一次元の$K_i$と一致する。
その非零列は$K_i$を張るので、
$K_i$が分解可能テンソルを含まないことは、
いずれかの列をreshapeした$2\times2$行列の
行列式が非零となることと同値である。
これがいずれかの$h_{ij}\ne0$という条件になる。

基底変更で$K_i$は
$(\operatorname{GL}(V_i)\otimes\operatorname{GL}(V_{i+1}))^\vee$
により移り、Segreの分解可能座も不変だから、
判定条件は基底によらない。
四つの$i$について適用すれば結論を得る。
\end{proof}

\subsection{既存の normal foundation chart の位置}

\begin{proposition}[一方の積を可逆にした表示]
\label{prop:qa-existing-foundation-chart}
\(b\) が可逆な開部分では \(F\simeq W\otimes V\) と同定でき、
結合律は \(d\) を \(a,c\) から一意に決める。
従って、この開部分の線形データは既存の normal foundation 表示
\[
 F_{\mathrm{old}}:P\otimes W_1\to W_2,
 \qquad G_{\mathrm{old}}:Q\otimes W_2\to W_3
\]
において
\((P,W_1,W_2,Q,W_3)=(V,U,E,W,N)\)、
\((F_{\mathrm{old}},G_{\mathrm{old}})=(a,c)\) としたものに一致する。
この同定の下で \((c,d)\) の全射性は \(c\) の全射性に等しい。
既存の表示の全開部分と比較する際には、
本節の分解可能な核を持たない条件を追加する。
\end{proposition}
\begin{proof}
\eqref{eq:qa-associativity} の右辺で \(b\otimes1_U\) が同型なので、
\(d=c(1_W\otimes a)(b^{-1}\otimes1_U)\) である。
残る空間と二写像は表示したものに一致する。
\(d\) は \(c\) を経由するので、両者の像の和は \(c\) の像である。
\end{proof}

この比較により、\cite[\texttt{prop:foundation-stack}]{QuadricModuli2026} に現れる一方の積を恒等化した chart と、
本節の二つの積がともに階数三になる境界とを区別できる。
古い二写像表示で \(\rank F_{\mathrm{old}}=2\) として得られる十二本矢表示は、
本節の許容開部分には入らない。
その核は二次元であり分解可能テンソルを含むからである。
また、その十二本矢は一方の対角位置に二組の追加矢を持つのに対し、
本節の最小十二本矢は両対角位置に一組ずつの追加矢を持つ。
矢の総数だけから、これらを同じ型と扱ってはならない。

% !TeX root = ../nc-quadric-surface-avatar.tex
\section{反標準 avatar の族と二つの十本矢 chart}
\label{sec:qa-avatar-charts}

本節では幾何的な族を直接定義し、その間の開埋め込みを証明する。
代数の同型類から幾何的 avatar への対応は
\cite[定義15.2・定理15.3]{Q12Avatar2026}に記載されているが、
その同型類分類を相対 AS 代数の stack 同値として使用しない。
族上の再埋め込みは\cite[第4--5節]{Q12Weyl2026}との比較も踏まえ、
本稿の補題~\ref{lem:qa-core-curve-vanishing}、
\ref{lem:qa-core-two-pencils}と定理~\ref{thm:qa-core-avatar-square-rotated}で
必要な幾何を自足的に証明し、
開部分の判定に必要なコホモロジー、非被約因子、垂直成分および降下を以下で確認する。
統合ノート\cite[第7.11--7.12項・第9節・第13.3項]{Synthesis2026}は
これらの出典と主張の水準を整理した資料であり、独立の検証の代用とはしない。

\paragraph{AS 側の辞書の対象.}
入力する体上分類の対象は
\cite[第15.1項・定理15.3]{Q12Avatar2026} の $A_Q^{\mathrm{pres}}$ である。
すなわち、局所有限な連結 positively directed $\mathbb Z$-代数で、
大域次元三、Gorenstein parameter 四の AS 正則性を満たし、
\eqref{eq:qa-universal-potential} の十二本矢 potential
（可縮な二次 pair を許す）の展開として表示できるものを、
全整数頂点を固定する代数同型で分類する。
最小化した四型では、右射影 $\mathsf P_i=e_iA$ と右単純 $\mathsf S_i$ に対して
\[
\begin{aligned}
0\longrightarrow\mathsf P_{i-4}
 &\longrightarrow\mathsf P_{i-3}^{\oplus2}\oplus\mathsf P_{i-2}^{\oplus\epsilon_i}\\
 &\longrightarrow\mathsf P_{i-1}^{\oplus2}\oplus\mathsf P_{i-2}^{\oplus\epsilon_i}
 \longrightarrow\mathsf P_i\longrightarrow\mathsf S_i\longrightarrow0
\end{aligned}
\]
という実際に完全な最小分解がある。
$\epsilon_i=0$、奇数 $i$ のみで $\epsilon_i=1$、
偶数 $i$ のみで $\epsilon_i=1$、全 $i$ で $\epsilon_i=1$ が、
順に八本、奇数対角十本、偶数対角十本、十二本に対応する。
右側の AS--Gorenstein 条件は
$\Ext_A^q(\mathsf S_i,\mathsf P_j)=0$ for $(q,j)\ne(3,i-4)$、
$\Ext_A^3(\mathsf S_i,\mathsf P_{i-4})\simeq\kk$ である。
分解の Euler 特性だけを課す条件ではない。
Noether 性はこの定義の追加仮定ではなく、引用する分類の結論である。
potential 表示、trivial pair の係数、周期同型、分解の基底は選択データにしない
\cite[定義2.1・第15.1項]{Q12Avatar2026}\cite[定義48.1]{AvatarVSeven2026}。
この整数頂点固定の同型は、四頂点 graded 代数の同型を射とした
$\mathfrak M_r^{\mathrm{alg}}$ と同じ相対同型層であるとは仮定しない。
本節の幾何的 stack の同型との比較も、体上の同型類辞書と区別する。

\subsection{分類する幾何データ}

基礎体 $\kk$ は標数零の代数閉体とする。
射影束は一次元商を分類する規約を用いる。
次の二つの標識付き曲面を考える。
\[
\begin{array}{c|c|c}
Y&h&f\\\hline
\Sigma _0=\PP^1\times\PP^1&\OO(1,1)&\OO(1,0)\\
\Sigma _2=\PP_{\PP^1}(\OO\oplus\OO(2))&\OO_{\PP(\OO\oplus\OO(2))}(1)&\pi^*\OO_{\PP^1}(1)
\end{array}
\]
いずれも $h^2=2$, $hf=1$, $f^2=0$, $-K_Y=2h$ である。
$\Sigma _2$ の負断面は $S=h-2f$ である。
本節の $\Pic^0$ は Picard 群の単位成分を表す。

\begin{definition}[曲線型の相対 avatar]
\label{def:qa-avatar}
局所 Noetherian な $\kk$-scheme $T$ 上の
$\mathcal A_{12}^{\times}$ の対象を
\[
 (Y\xrightarrow{\pi}B\longrightarrow T;\ h,f;\ D\hookrightarrow Y;\ P)
\]
とする。ここで $B/T$ は $\PP^1$-fibration、$Y/B$ は
$T$ 上 fppf 局所で $\PP_B(\mathcal E)$ と表示される射影直線族であり、
各幾何 fiber の $(Y,h,f)$ は上表のいずれかである。
$h,f$ は相対 Picard の類で、$f$ は指定した ruling に対応する。
$D$ は $T$ 上平坦な相対有効 Cartier 因子で、相対因子類は $2h$ である。
$P$ は $\Pic^0_{D/T}$ の section で、各幾何 fiber 上で非自明とする。
射は ruling、標識、$D$ および $P$ の類を保つ $T$ 上の同型である。
線束の scalar trivialization はデータに加えない。
\end{definition}

この定義は $\Sigma _0$ と $\Sigma _2$ の間で型が跳ぶ族を許す。
実際の線束を用いる式は fppf 局所で解釈し、基底からの線束の tensor は区別しない。
以下で
\[
 p=\pi|_D:D\longrightarrow B,\qquad F=\OO_Y(f)|_D,\qquad H=\OO_Y(h)|_D
\]
と書く。特に $H$ は線束であり、内在的モジュライの記号とは無関係である。
曲面、因子および相対 Picard 類の降下により、定義\ref{def:qa-avatar}は
fppf stack を定める。本節の「開部分 stack」は任意の基底への引戻しが
開部分 scheme となる意味で用いる。

\begin{remark}[有限表示と scalar の扱い]
fppf 局所で $B=\PP^1_T$ とし、$\mathcal E=\pi_*\OO_Y(h)$ を取ると、
各幾何 fiber の分裂型は $\OO(1)^{\oplus2}$ または $\OO\oplus\OO(2)$ である。
評価列は
\[
 0\longrightarrow \mathcal R\otimes\OO_B(-1)
 \longrightarrow \mathcal V\otimes\OO_B
 \longrightarrow\mathcal E\longrightarrow0,
 \qquad\rank\mathcal R=2,\quad\rank\mathcal V=4
\]
となる。したがって線形行列、その余核の射影束上の相対反標準因子、
および相対 Picard 類によって局所的な parameter を記述できる。
この表示の基底変更をすべて新たな自己同型として保持すると、
定義\ref{def:qa-avatar}と異なる stack になる。
本節の射はあくまで標識付き曲面と因子の同型である。
\end{remark}

\subsection{非被約因子を含むコホモロジー}

\begin{lemma}[支持に対する判定]
\label{lem:qa-support}
体上で $D\in|2h|$ とする。このとき
$\omega_D\simeq\OO_D$、$h^0(\OO_D)=h^1(\OO_D)=1$ である。
真の非零部分因子 $0<Z<D$ について $Z(D-Z)\geq2$ が成り立つ。
線束 $N$ の非零切断が正確に $Z\leq D$ に支持されるなら
\[
 \deg_ZN\ \geq\ Z(D-Z).
\]
従って、線束 $L$ がすべての $0<Z\leq D$ について
\begin{equation}
 \deg_ZL+Z(D-Z)>0
 \label{eq:qa-support-vanishing}
\end{equation}
を満たせば $H^1(D,L)=0$ である。
\end{lemma}

\begin{proof}
随伴と $0\to\omega_Y\to\OO_Y\to\OO_D\to0$ により
双対化層と構造層のコホモロジーを得る。
$\Sigma _0$ では $D=2f+2g$、$Z=af+bg$ と書けて
\[
 Z(D-Z)=2(a+b-ab),\qquad 0\leq a,b\leq2.
\]
$Z=0,D$ を除いたすべての場合で右辺は二以上である。
$\Sigma _2$ では $D=2S+4f$、$Z=aS+bf$ と書くと
\[
 Z(D-Z)=2a^2+2b(1-a),\qquad0\leq a\leq2,\quad0\leq b\leq4.
\]
$a=0$ では $b\geq1$、$a=1$ では右辺が二、
$a=2$ では真部分因子なら $b\leq3$ なので、同じ結論を得る。

切断の generic annihilator によって、その正確な支持である部分因子 $Z$ を取る。
各既約成分の generic な Artin 局所環で切断の消滅次数を取り除いて $Z$ の重複度を定めると、
切断は $N|_Z(-(D-Z))$ の、各 generic point で可逆な切断になる。
$Z$ は Cohen--Macaulay なので、この切断が定める
$\OO_Z\to N|_Z(-(D-Z))$ は単射で、余核は有限長である。
Euler 特性を比較すれば不等式を得る。
これは部分因子を重複度込みで取る議論であり、$D$ の被約性を使わない。
最後は $\omega_D\simeq\OO_D$ による Serre 双対を用い、
$L^{-1}$ の切断にこの不等式を適用すればよい。
\end{proof}

\begin{lemma}[再埋め込みに必要な消滅]
\label{lem:qa-vanishing}
$R\in\Pic^0(D)$ とする。このとき
\[
 H^1(D,HR)=H^1(D,HF^{-1}R)=0,\qquad
 \chi(HR)=4,\quad\chi(HF^{-1}R)=2.
\]
また $R\not\simeq\OO_D$ なら $H^*(D,R)=0$ である。
\end{lemma}

\begin{proof}
$R$ は各部分因子上で次数零である。
$h$ は nef で全次数が四なので、補題\ref{lem:qa-support}を適用できる。
$\Sigma _0$ では $h-f$ も nef で全次数が二である。
$\Sigma _2$ では $Z=aS+bf$ 上で
\[
 \deg_Z(h-f)+Z(D-Z)=2a^2-a+3b-2ab.
\]
$a=0,1,2$ に対する右辺はそれぞれ $3b$, $1+b$, $6-b$ であり、
$Z\neq0$ の許される範囲で正である。Riemann--Roch が Euler 特性を与える。
$R$ の非零切断は真部分因子には支持されない。
全体に支持される切断の余核は次数零なので零であり、$R$ は自明になる。
$R^{-1}$ にも同じ議論を適用し、Serre 双対を使う。
\end{proof}

\subsection{Picard twist による相対的な再埋め込み}

\begin{proposition}[垂直成分がある場合も含む直像]
\label{prop:qa-pushforward}
定義\ref{def:qa-avatar}の族で、$R\in\Pic^0_{D/T}(T)$ を取る。
fppf 局所で線束によって表示すると
\[
 Rp_*(HR)\simeq\mathcal E_R:=p_*(HR)
\]
は $B$ 上の階数二の局所自由層であり、任意の基底変換と両立する。
各幾何 fiber $B_t\simeq\PP^1$ 上で次数二かつ大域生成である。
評価写像は閉埋め込み
\[
 j_R:D\hookrightarrow Y_R:=\PP_B(\mathcal E_R)
\]
を定め、像は相対反標準 Cartier 因子である。
新しい標識の制限は $F_R=F$、$H_R=HR$ である。
\end{proposition}

\begin{proof}
まず体上で $b\in B$、$G_b=\pi^{-1}(b)$ とする。
$G_b$ が $D$ の成分でなければ $D\cap G_b$ は長さ二の scheme であり、
導来 fiber の切断は二次元で次数零に集中する。
$G_b\leq D$ なら、局所 parameter $t$ と $D$ の方程式 $t^m a$、$t\nmid a$ を用いると
\[
 \operatorname{ann}_{\OO_Y/(t^ma)}(t)=(t^{m-1}a)/(t^ma).
\]
したがって $L=HR$ の $B$ 方向の導来制限は
\[
 \mathcal H^{-1}=L|_{G_b}\otimes\OO_{G_b}(-D+G_b)\simeq\OO_{\PP^1}(-1),
 \qquad
 \mathcal H^0=L|_{G_b}\simeq\OO_{\PP^1}(1)
\]
を持つ。ここで $\deg_{G_b}L=1$、$D G_b=2$、$G_b^2=0$ を使った。
次数 $-1$ の層は acyclic であるから、この場合も直像の導来 fiber は $\kk(b)^2[0]$ となる。
垂直成分を持つ $p$ は有限でも平坦でもないが、この計算にはその仮定を要しない。

$i_*L$ は Cartier immersion $i:D\hookrightarrow Y$ の線束の直像なので perfect である。
滑らかな射影的射 $\pi$ による perfect pushforward と導来基底変換を用いると、
$R\pi_*i_*L$ は perfect である。
すべての導来 fiber が同じ階数二の次数零複体なので、局所的に可逆 minor で
可縮項を消去することにより、階数二の局所自由層となる。
族の場合、$i_*L$ は $T$ 上平坦であり、まず $T$ の幾何 fiber、次に $B_t$ の一点へ
制限すれば同じ計算になる。このため非被約な基底も含めて結論が成り立つ。

補題\ref{lem:qa-vanishing}と Riemann--Roch により
\[
 \deg\mathcal E_R=\chi(HR)-2=2,\qquad
 H^1(B,\mathcal E_R(-1))=H^1(D,HF^{-1}R)=0.
\]
$\PP^1$ 上の分裂定理より、分裂型は $\OO(1)^{\oplus2}$ または $\OO\oplus\OO(2)$ である。

各 $B$-fiber での評価は、長さ二の scheme 上の完全切断の評価、
または $\OO_{\PP^1}(1)$ の評価である。従って全射で、誘導される fiber の射は閉埋め込みである。
Nakayama の補題から全評価も全射となる。$j_R$ は proper かつ quasi-finite なので有限であり、
有限射の基底変換と再度の Nakayama により
$\OO_{Y_R}\to(j_R)_*\OO_D$ は全射である。
体上では像は滑らかな曲面内の Cohen--Macaulay な余次元一部分 scheme なので Cartier である。
像の類を $a h_R+b f_R$ とすると、ruling との交点から $a=2$、
$h_R$ との交点 $\deg_D(HR)=4$ から $b=0$ を得る。
族では $D/T$ の平坦性と各 fiber の Cartier 性によって相対 Cartier 性が従う。
各 fiber で自明な差の線束は、$H^1(\OO_{Y_t})=0$ と基底変換により基底から来るので、
相対 Picard 類も $2h_R$ である。
\end{proof}

\begin{theorem}[族上の幾何的並進]
\label{thm:qa-translation}
前命題において $R=P^n$ とした構成は $\mathcal A_{12}^{\times}$ の自己同値
$\mathsf T^n$ を定める。自然な同型
\[
 \mathsf T^m\mathsf T^n\simeq\mathsf T^{m+n},\qquad
 \mathsf T^0\simeq\operatorname{id}
\]
があり、逆操作は $P^{-n}$ を用いる同じ構成である。
\end{theorem}

\begin{proof}
抽象的な $D$、ruling $p:D\to B$ および $P$ は変更せず、$H$ を $HP^n$ に変える。
二回適用した直像は $p_*(HP^nP^m)$ なので、tensor の結合則が合成則を与える。
$R=\OO_D$ の場合、制限列
\[
 0\longrightarrow\OO_Y(-h)\longrightarrow\OO_Y(h)\longrightarrow i_*H\longrightarrow0
\]
と $R\pi_*\OO_Y(-h)=0$ が元の $\mathcal E=\pi_*\OO_Y(h)$ を回収する。
相対的には左項に基底線束の tensor が入ることがあるが、同じ消滅が成り立つ。
評価の自然性により射影束だけでなく $D$ の埋め込みも回収する。
線束の代表を基底線束で tensor すると直像も同じ線束で tensor されるため、
射影束および相対標識は変わらない。
$q:D\to T$ について、各幾何 fiber で $h^0(\OO_D)=1$ であるから、
コホモロジーと基底変換により $q_*\OO_D=\OO_T$ が任意の基底変換後も成立する。
従って線束代表の自己同型は基底からの scalar であり、射影化で消える。
従って構成と合成同型は fppf 降下し、射と基底変換に整合する。
\end{proof}

この定理は\cite[定理5.4]{Q12Weyl2026}の操作である。
ここで証明したのは幾何的 avatar の stack 上の並進であり、
無限二面体群のすべての関係を stack 上へ持ち上げる主張ではない。

\subsection{二つの開 chart と四つの層}

\begin{definition}[幾何的な八本矢・十本矢 chart]
\label{def:qa-open-charts}
$\mathcal U\subset\mathcal A_{12}^{\times}$ を、全幾何 fiber が $\Sigma _0$ である部分とする。
\[
 \mathcal A_{10,\mathrm e}^{\times}:=\mathcal U,
 \qquad
 \mathcal A_{10,\mathrm o}^{\times}:=\mathsf T^{-1}(\mathcal U),
 \qquad
 \mathcal A_8^{\times}:=\mathcal U\cap\mathsf T^{-1}(\mathcal U)
\]
と置く。添字 $\mathrm e,\mathrm o$ は、後述する標準表示で残る偶数対角・奇数対角に対応する。
ここでは定義自体を幾何データだけで与えている。
\end{definition}

\begin{theorem}[avatar の開部分 stack]
\label{thm:qa-open-charts}
次の射は開埋め込みである。
\[
 \mathcal A_8^{\times}\hookrightarrow
 \mathcal A_{10,\mathrm o}^{\times}\hookrightarrow\mathcal A_{12}^{\times},
 \qquad
 \mathcal A_8^{\times}\hookrightarrow
 \mathcal A_{10,\mathrm e}^{\times}\hookrightarrow\mathcal A_{12}^{\times}.
\]
幾何点での層は次の表で与えられる。
\begin{center}
\begin{tabular}{c|c|c}
曲面&曲線型の条件&属する層\\\hline
$\Sigma _0$&$P\not\simeq F^2H^{-1}$&$\mathcal A_8^{\times}$\\
$\Sigma _0$&$P\simeq F^2H^{-1}$&$\mathcal A_{10,\mathrm e}^{\times}\setminus\mathcal A_8^{\times}$\\
$\Sigma _2$&$D\cap S=\varnothing$&$\mathcal A_{10,\mathrm o}^{\times}\setminus\mathcal A_8^{\times}$\\
$\Sigma _2$&$S\leq D$&$\mathcal A_{12}^{\times}\setminus
(\mathcal A_{10,\mathrm e}^{\times}\cup\mathcal A_{10,\mathrm o}^{\times})$
\end{tabular}
\end{center}
最後の行は幾何点集合として閉 locus である。
全 fiber がこの型となる族の充満部分関手に、閉埋め込みを自動的に与える主張ではない。
\end{theorem}

\begin{proof}
$\mathcal E=\pi_*\OO_Y(h)$ の二つの分裂型は
\[
 h^0(B_t,\mathcal E_t(-2))=
 \begin{cases}0&Y_t=\Sigma _0,\\1&Y_t=\Sigma _2\end{cases}
\]
によって区別される。コホモロジーの上半連続性から $\mathcal U$ は開である。
自己同値 $\mathsf T$ による逆像と交わりも開なので、すべての開埋め込みを得る。
残るのは幾何点での条件である。

$\Sigma _0$ で $H=FG$ と書くと
\[
 h^0(B,\mathcal E_P(-2))=h^0(D,GF^{-1}P).
\]
$Z=af+bg$ に対し、$GF^{-1}P$ の次数は $a-b$ である。
真部分因子では
\[
 Z(D-Z)-\deg_Z(GF^{-1}P)=a+3b-2ab>0.
\]
よって補題\ref{lem:qa-support}により、この線束の非零切断は全体に支持される。
全次数は零なので、非零切断があることは線束が自明であることと同値である。
従って $\mathsf T(Y)=\Sigma _2$ となる条件はちょうど $P\simeq F^2H^{-1}$ である。

$\Sigma _2$ で $D\cap S=\varnothing$ なら $H\simeq F^2$ だから
$h^0(D,HF^{-2}P)=h^0(D,P)=0$ であり、$\mathsf T(Y)=\Sigma _0$ となる。
$S\leq D$ なら $C=D-S$ と置く。$C\sim h+2f$ であり、
$\OO_Y(-C)$ は $\pi$ に沿って次数 $-1$ なので acyclic である。
従って $h^0(\OO_C)=1$, $h^1(\OO_C)=0$、標数零より $\Pic^0(C)=0$ を得る。
残余因子の完全列は
\[
 0\longrightarrow P|_C\longrightarrow\OO_D(S)P
 \longrightarrow\OO_S(-2)\longrightarrow0.
\]
$P|_C\simeq\OO_C$ なので $h^0(D,HF^{-2}P)=1$ である。
従って $\mathsf T(Y)=\Sigma _2$ のままである。
最後に $D S=0$ であるから、有効因子 $D$ は $S$ を含むか $S$ を避けるかのいずれかである。
これですべての場合を得る。閉性は二つの開部分の和の補集合であることによる。
\end{proof}

\begin{remark}[矢数の名称と体上の辞書]
幾何的な開埋め込みと表の分解は、上の証明の通り
射影束の分裂型、Picard 類および垂直成分条件から従う。
各層を $8,10,10,12$ 本の最小矢数と対応させる体上の解釈は、
$A_Q^{\mathrm{pres}}$ を対象とする
\cite[定理15.3]{Q12Avatar2026}の同型類分類を用いる。
同じ矢数の相対 QP 表示を比較する際の階数条件は、
上の幾何学的な層の名称からは
推論せず、各代数表示の積と関係について検証する。
\end{remark}


\begin{remark}[「十本矢全体」の二通りの意味]
最小十本矢の二つの層の和は $\mathcal A_8^{\times}$ を含まない。
これに対し上の $\mathcal A_{10,\mathrm o}^{\times}$ と
$\mathcal A_{10,\mathrm e}^{\times}$ は、可縮な pair を許す十本矢表示の
幾何的 chart に対応し、八本矢部分を開部分として含む。
二つの十本矢 chart の和は、最深の十二本矢層の補集合である。
この区別なしに「矢の少ないモジュライが矢の多いモジュライの開部分である」
と述べることはできない。
\end{remark}

\subsection{直積型 collection と全 square avatar}

$\mathcal A_{\mathrm{sq}}^{\times}$ を、順序付きの二つの ruling を持つ
$\Sigma _0$-族、その相対 $(2,2)$ Cartier 因子 $D$、および fiberwise 非自明な
$P\in\Pic^0_{D/T}$ の stack とする。
曲線上の ruling の制限を $F_s,G_s$ と書く。
これは $h=f+g$ と置くことによって $\mathcal U$ と同じ幾何的 stack である。

\begin{corollary}[全 square chart の相対的な対応]
\label{cor:qa-square-chart}
幾何的 stack の同値
\[
 \mathcal A_{\mathrm{sq}}^{\times}
 \xrightarrow[\ \mathsf T^{-1}\ ]{\ \sim\ }
 \mathcal A_{10,\mathrm o}^{\times}
\]
があり、制限線束は
\begin{equation}
 F=F_s,\qquad H=F_sG_sP^{-1},\qquad P=P
 \label{eq:qa-square-dictionary}
\end{equation}
となる。逆には $\mathsf T$ によって得られる $\Sigma _0$ 上で
\[
 F_s=F,\qquad G_s=HPF^{-1}
\]
を取る。square 座標で八本矢部分となる条件は
$P\not\simeq G_sF_s^{-1}$ であり、壁 $P\simeq G_sF_s^{-1}$ は
負断面を避ける $\Sigma _2$ 上の奇数対角十本矢層に移る。
\end{corollary}

\begin{proof}
定義から $\mathsf T^{-1}:\mathcal U\to\mathsf T^{-1}(\mathcal U)$ は同値であり、
その線束の式が\eqref{eq:qa-square-dictionary}である。
逆の式は $H_s=HP=F_sG_s$ による。
元の $\Sigma _0$ で $\mathsf T^{-1}$ が $\Sigma _2$ を与える条件は、
定理\ref{thm:qa-open-charts}の計算を $P^{-1}$ に対して行って
$P^{-1}=F_sG_s^{-1}$、すなわち $P=G_sF_s^{-1}$ である。
その像は $\mathcal A_{10,\mathrm o}^{\times}$ にあるので負断面を避ける。
\end{proof}

この系では反標準因子を滑らか、被約または既約と仮定していない。
従って「十本でちょうど全ての quadric の divisor と $\Pic^0$ が出る」
という記述は、\emph{奇数対角の非最小十本矢 chart と全 square avatar の対応}
として正確に定式化できる。最小十本矢だけの主張ではなく、$P=\OO_D$ も除いている。

代数側の square foundation は
\[
 B_{\mathrm{sq}}=\End_{\to}(P,F_s,G_s,F_sG_s)
\]
であり、可換な基準 collection は
$(\OO,\OO(1,0),\OO(0,1),\OO(1,1))$ である
\cite[定義1.2・定理1.3]{AvatarVSeven2026}。
最初の二項の左 mutation は曲線上で
\[
 (P^2F_s^{-1},P,G_s,F_sG_s)
\]
を与える。共通に $F_sP^{-1}$ を tensor すると
\[
 (P,F_s,F_sG_sP^{-1},F_s^2G_sP^{-1})=(P,F,H,HF)
\]
となり、上の幾何的辞書と一致する。
mutationによる指定型AS代数の対応は定理~
\ref{thm:qa-as10-global-bijection}と
\ref{thm:qa-as10-relative-families}で証明する。
元の体上および仮定付き相対対応は
\cite[定理51.4・系51.5、定理53.2]{AvatarVSeven2026}にある。
幾何的な開埋め込みは独立に証明した。

\subsection{滑らかな楕円型における三者の明示的辞書}
\label{subsec:qa-square-elliptic-pointed-dictionary}

\begin{theorem}[VdB楕円四つ組とsquare avatarの同型類の全単射]
\label{thm:qa-elliptic-vdb-square-avatar-bijection}
$k$を標数零の代数閉体とする。
滑らかな楕円曲線$E$上のVan den Bergh許容四つ組
$(E,L_0,L_1,L_2)$で、各$L_i$は次数二、
\[
 L_1L_0^{-1}\ne\OO_E,\qquad
 L_2L_0^{-1}\ne\OO_E,\qquad
 L_1L_2^{-1}\ne\OO_E
 \tag{VdB-adm}\label{eq:qa-vdb-square-admissibility}
\]
を満たすものの、曲線・三線束を保つ同型類を考える。
これは、\emph{順序付き}二rulingを持つ
滑らかなsquare avatar
$(Y=\PP^1\times\PP^1,D,P)$で
\[
 D\in|\OO_Y(2,2)|,\qquad D\text{ は滑らか},\qquad
 P\in\Pic^0(D)\setminus\{\OO_D\},\qquad
 P\not\simeq G_sF_s^{-1}
 \tag{Sq-cub}\label{eq:qa-square-elliptic-cubic-open}
\]
を満たすものの同型類と自然に全単射する。
ここで$F_s=\OO_D(1,0)$、
$G_s=\OO_D(0,1)$とし、
二rulingの順序は同型の中で保持する。

全単射の\emph{両方向の式}は
\[
 \boxed{
 \begin{aligned}
 (E;L_0,L_1,L_2)&\longmapsto
  (D\hookrightarrow\PP H^0(F_s)^\vee
           \times\PP H^0(G_s)^\vee;\ P),\\
 F_s&=L_2,\qquad
 G_s=L_1L_2L_0^{-1},\qquad
 P=L_2L_0^{-1},\\
 (D;F_s,G_s,P)&\longmapsto
  (E=D;\ L_0=F_sP^{-1},\
   L_1=G_sP^{-1},\ L_2=F_s).
 \end{aligned}}
 \tag{Sq-VdB}\label{eq:qa-square-vdb-explicit}
\]
で与えられる。
この同型類の比較は、周期同型や根を
追加標識にしないVdBのcubic$\mathbb Z$-代数
の幾何的分類を通じて、
\emph{選んだcubic normal foundation}の
相対的な$k$点の辞書になる。

差パラメータ
$\eta=L_1L_0^{-1}$、
$\delta=L_2L_0^{-1}$を用いれば、
\[
 \eta=G_sF_s^{-1},\qquad \delta=P.
 \tag{Sq-diff}\label{eq:qa-square-vdb-differences}
\]
\end{theorem}
\begin{proof}
まず$L_0,L_1,L_2$から
$P=L_2L_0^{-1}$、
$F_s=L_2$、
$G_s=L_1L_2L_0^{-1}$を作ると、
$F_s,G_s$はいずれも次数二で、
$G_sF_s^{-1}=L_1L_0^{-1}$である。
式~\eqref{eq:qa-vdb-square-admissibility}より
$F_s\not\simeq G_s$である。
滑らかな楕円曲線上の相異なる次数二線束の
完全線形系は、それぞれ二重被覆
$E\to\PP^1$を与え、その積は
$\PP^1\times\PP^1$への閉埋込みである。
像は$(2,2)$因子で、反標準因子となる。
また$P\ne\OO_E$、
$P\not\simeq G_sF_s^{-1}$である。
これは曲線同型を保ち、線形系の基底の選択を
$\PGL_2^2$で割ればavatarの同型類を定める。

逆にsquare avatarから
$L_0=F_sP^{-1}$、
$L_1=G_sP^{-1}$、
$L_2=F_s$を作ると、
式~\eqref{eq:qa-square-elliptic-cubic-open}により
$\eta=G_sF_s^{-1}\ne\OO_E$、
$\delta=P\ne\OO_E$、
$\eta\ne\delta$となる。
従ってVdBの許容条件が成立する。
二つの式を順次適用すると三線束と
二つのrulingの完全線形系をそのまま回収するので、
互いに逆である。
VdBの幾何的quintuple復元
\cite[Theorem 5.5.10]{VdB2011}を適用すれば
最後の解釈も従う。
\end{proof}

\begin{theorem}[O-markingを保つruling反射と楕円型の商全単射]
\label{thm:qa-elliptic-avatar-pointed-quotient}
前定理のsquare avatarについてさらに
\[
 G_sF_s^{-1}\not\simeq P^2
 \tag{two-cub}\label{eq:qa-square-two-cubic-open}
\]
を課した開を$\mathcal A_{\mathrm{sq,ell}}^{2c}$とする。
この開上では
\[
 \boxed{\quad
 r_1(F_s,G_s,P)
  =(G_sP^{-1},F_sP,P)
 \quad}
 \tag{Sq-R}\label{eq:qa-square-pointed-reflection}
\]
は、右辺の二つの次数二線束が定める
新しい$(2,2)$埋込みを用いて、
幾何的avatarの同型類の対合を定める。
前定理の式~\eqref{eq:qa-square-vdb-explicit}の下で、
これは論文Bの
\cite[\texttt{lem:elliptic-two-reflections}]{QuadricModuli2026}
の\emph{構造対象$\OO$を保つ}ruling交換
\[
 (L_0,L_1,L_2)\longmapsto
 (L_0L_1L_2^{-1},L_2,L_1)
\]
に一致する。
従って、$y=(\mathcal C,\OO)$が滑らかな
楕円型かつ二つのnormal foundationがともに
cubic型となる\emph{幾何的なO-marked同型類集合}は
\[
 \boxed{\quad
 \bigl|\mathcal A_{\mathrm{sq,ell}}^{2c}(k)\bigr|
    /\langle r_1\rangle
 \ \xrightarrow{\ \sim\ }\
 \bigl|\mathfrak H_{2,\mathrm{ell}}^{2c}(k)\bigr|
 \quad}
 \tag{Sq-H2}\label{eq:qa-square-pointed-bijection}
\]
と全単射する。
左辺はavatar同型類の集合をとってから
$r_1$の軌道で割ることを意味する。
ここでは$S_2$の商の全相対stack同値を
この幾何的全単射だけから推論しない。

一方、helixの一段の移動に対応する
\[
 r_0(F_s,G_s,P)=(G_s,F_sP,P)
 \tag{Sq-T}\label{eq:qa-square-unpointed-shift}
\]
はVdB四つ組の
$(L_0,L_1,L_2)\mapsto(L_1,L_2,L_3)$に
対応するが、一般に$\OO$を変更する。
従って\emph{O-marked}同型類を得るために
$r_0$まで同一視してはならない。

この商の記述は、一般の二つのavatarの
$r_1$-軌道が二元であることを要求しない。
特殊な二捩れ慣性点では二つが同型になる。
\end{theorem}
\begin{proof}
式~\eqref{eq:qa-square-vdb-explicit}を
$r_1$の右辺へ適用すると
\[
 \begin{aligned}
 L'_0&=G_sP^{-2}
       =L_0L_1L_2^{-1},\\
 L'_1&=F_s=L_2,\qquad
 L'_2=G_sP^{-1}=L_1
 \end{aligned}
\]
を得る。新しい二ruling線束が相異なる条件は
$G_sP^{-1}\not\simeq F_sP$、
すなわち式~\eqref{eq:qa-square-two-cubic-open}
に他ならない。
もう一度適用すると$(F_s,G_s,P)$を回収するため
対合である。

論文Bの二つの内在的foundationの系
\cite[\texttt{cor:two-intrinsic-foundations}]{QuadricModuli2026}と、
pointed groupoid比較の系
\cite[\texttt{cor:h2-point-groupoid-foundation}]{QuadricModuli2026}は、指定$\OO$から始まるnormal foundationは
exact同値の下でちょうど二択であり、
二つともcubic型なら他方は上記$R$で得られること、
またfoundation同型はpointed exact同型を
完全に検出することを示す。
Van den Berghの四つ組復元を合わせれば、
二cubic楕円型のpointed同型類に関する
核と像はちょうど$r_1$の同型類軌道となる。
これが式~\eqref{eq:qa-square-pointed-bijection}である。

同様に$r_0$の右辺から得る
$(L'_0,L'_1,L'_2)$は
\[
 (G_sP^{-1},F_s,G_s)=(L_1,L_2,L_3),
 \quad L_3=L_0^{-1}L_1L_2
\]
である。論文Bの同じ補題では、この添字移動は
underlying Abel圏の同値類は保つが、
一般に構造対象$\OO$を保たない。
従って$r_0$を割るのは別の
unpointed moduli問題である。
\end{proof}

\begin{corollary}[OU potentialからsquare avatar、O-marked対象への一般次数]
\label{cor:qa-ou-potential-square-avatar-pointed-degrees}
論文C
~\ref{thm:qa-ou-rigid-potential-quintuple-generic-etale}
の正則半単純な滑らかな楕円開を、
前定理の二cubic開~\eqref{eq:qa-square-two-cubic-open}
および一般慣性同型の稠密開へ制限する。
この開ではpotentialの幾何的四重表示が
VdB四つ組$(E,L_0,L_1,L_2)$を与え、
式~\eqref{eq:qa-square-vdb-explicit}により
\emph{square avatarへの相対的な対応}がある。
全ての線束・曲線が相対的に構成される
étale chart上では、これは基底変更と可換な
代数スタックの射となり、
相対的な曲線同型を通じて降下する。
一方、対応の幾何的fiberの一般次数は
\[
 \boxed{
 \begin{array}{ccc}
 \text{有効potential表示}
   &\xrightarrow{\;4\;}&
  \text{順序付きsquare avatar}\\
 &&\big\downarrow\vcenter{\rlap{$\scriptstyle 2$}}\\
 &&\text{O-marked }\mathfrak H_2
 \end{array}}
 \tag{Pot-Sq-H2}\label{eq:qa-potential-square-pointed-degrees}
\]
であり、potentialからO-marked対象への
一般次数は$4\cdot2=8$である。
ここで上の$4$はpotential/quintupleの
\emph{表示間}の一般四重射、
下の$2$は二つのnormal foundationの忘却であり、
$\Pic^0$の四つの平方根の選択を
追加標識として数えた次数ではない。
\end{corollary}
\begin{proof}
Cの
~\ref{thm:qa-ou-rigid-potential-quintuple-generic-etale}
による一般有限étale四重射は、
指定された幾何quintupleの相対四重テンソルを与える。
滑らかな正則半単純開では、C
~\ref{cor:qa-ou-relative-regular-semisimple-potential}
の点曲線と二rulingの相対構成が、
基底変更と可換な三つの次数二線束を与える。
それらのテンソル積と二つの完全線形系で
式~\eqref{eq:qa-square-vdb-explicit}の
square avatarを相対的に作る。
曲線・線束の同型を保つので基底変更と降下に
関して自然である。
VdBの四つ組の幾何的同型類とsquare avatarの
全単射は定理
~\ref{thm:qa-elliptic-vdb-square-avatar-bijection}で証明した。
これとCの
~\ref{cor:qa-quintuple-h2-generic-degree-two}の
二つのfoundationの一般次数を合成すれば
表示された幾何一般次数が従う。
\end{proof}

\begin{remark}[一cubic境界と古典点での辞書の射程]
\label{rem:qa-square-avatar-boundary-cautions}
$G_sF_s^{-1}=P^2$では、$r_1$後の
$F'_s=G_sP^{-1}$と$G'_s=F_sP$が一致し、
square$(2,2)$埋込みを構成できない。
これは論文C
~\ref{cor:qa-ou-one-cubic-fiber-and-nonfiniteness}
の一cubic境界であり、
二つ目のnormal foundationはnoncubicとして
内在的な選択被覆
~\ref{thm:qa-global-two-ruling-etale-cover}
に残る。

また$P=\OO_D$を任意の$(2,2)$因子に
そのまま許すと、これらの曲線から得る
可換な四項foundation代数は$D$に依存せず、
同じ可換quadricを重複して数える。
従って全同型類の全単射には古典点を
別途一点として追加し、元のavatarの
$P=0$座を単に追加するだけでは足りない。
$\Sigma_2$の古典型は別のweighted
stackを生み、無条件には$\HHstack$に入らない。
$\Sigma_2$曲線型でのO-marked適合性、
非被約反標準因子を含む全対応、
およびpotentialの安定化・自己同型群の比較は
式~\eqref{eq:qa-square-pointed-bijection}
の証明には含まれない。
\end{remark}

\begin{proposition}[二つのwallは異なり、noncubic avatar壁には追加認識が必要]
\label{prop:qa-avatar-walls-and-noncubic-dimension}
滑らかな楕円型square avatarについて、
$\eta=G_sF_s^{-1}$、$\delta=P$と書く。
次の二つのwallを区別する。
\begin{enumerate}
\item $\eta=2\delta$、$\delta\notin E[2]$：
式~\eqref{eq:qa-square-vdb-explicit}の四つ組は
geometric cubic型であるが、
$r_1$後の第二foundationはnoncubicになる。
したがって元のO-marked点は
\emph{cubic locusに属する}。
\item $\eta=\delta\ne0$：
式~\eqref{eq:qa-square-vdb-explicit}の
$L_1=L_2$なので元のfoundationが
geometric cubic型ではない。
このwallはsquareの八本矢開の補集合であり、
曲線型avatarの十本矢・$\Sigma_2$ chartへ
移行する。
\end{enumerate}
後者の全smooth楕円avatar同型類の
パラメータ空間は、楕円曲線の$j$と
非零点$\delta\in E$を局所座標として
\emph{二次元}である。
一方、論文Bの内在noncubic境界
$\mathfrak N\subset\HHstack$は
\emph{純一次元}のArtinスタックである。
従ってこのwall\emph{全体}を
内在noncubic境界と表現可能quasi-finiteな
モジュライ同型で結び付けることはできない。
O-marked適合性はこのwall上で
さらに制限されるか、あるいは
同じpointed対象へ正の相対次元を持って
潰れるかのいずれかである。

特に非零二捩れ$P$の条件$P^2=\OO_E$は
楕円曲線のmoduli上で一次元を定め、
次に検討すべき\emph{候補}ではあるが、
それが内在noncubic境界をちょうど尽くすことは
ここでは証明していない。
\end{proposition}
\begin{proof}
最初のwallでは$L_0=F_sP^{-1}$、
$L_1=G_sP^{-1}$、$L_2=F_s$が
互いに相異なり、VdB許容条件を満たす。
$r_1$による新しい二線束は
$G_sP^{-1}=F_sP$と$F_sP$で一致するので、
第二はnoncubicとなる。
これはC
~\ref{cor:qa-ou-one-cubic-fiber-and-nonfiniteness}
の一cubic境界に一致する。
次のwallでは$G_s=F_sP$だから
$L_1=G_sP^{-1}=F_s=L_2$となる。
$\mathcal A_8^\times$の条件
~\eqref{eq:qa-square-elliptic-cubic-open}
に反することも直接に分かる。
$\Sigma_2$との対応は
系~\ref{cor:qa-square-chart}に従う。

最後の次元を計算する。
楕円曲線$E$のmoduliは一次元で、
$F_s\in\Pic^2(E)$の選択は曲線の
translationによって動かせる。
標数零では$[2]:E\to E$が全射なので、
このtranslationの軌道は$\Pic^2(E)$全体である。
$\eta=\delta$の下では$G_s=F_sP$だから
残るパラメータは
$E$と$0\ne P\in\Pic^0(E)$で、
そのmoduliの一般次元は$1+1=2$。
一方、論文B
\cite[\texttt{thm:noncubic-boundary-smooth-codim-two}]{QuadricModuli2026}
は全$\HHstack$が純次元三で
$\mathfrak N$が滑らかな余次元二、
すなわち純次元一であることを証明している。
有限型quasi-finite射の像は源の各成分と同じ
次元を持つので、この二次元wall全体を
一次元の境界にquasi-finiteに写すことはできない。
最後の二捩れ条件の次元も
有限étaleな$E[2]\setminus0$から従う。
\end{proof}

\begin{theorem}[二捩れsquare壁とnoncubic二根対の幾何同型類の全単射]
\label{thm:qa-avatar-two-torsion-root-pair-bijection}
$k$を標数零の代数閉体とする。
次の二集合を考える。
\begin{enumerate}
\item 滑らかな楕円曲線$E$、次数二線束$F_s$、
非自明な$P\in\Pic^0(E)[2]$の三つ組
$(E,F_s,P)$の曲線・線束・Picard類を
保つ同型類集合。$G_s=F_sP$とすると
$F_s\not\simeq G_s$だから、二完全線形系は
$(2,2)$埋込みを定め、対応するsquare avatarは
壁$P=G_sF_s^{-1}$上にある。
\item 射影直線$B\simeq\PP^1$上の、
互いに交わらない\emph{被約}二点因子の
非順序対$\{B_1,B_2\}$で、
$\deg B_i=2$かつ$B_1+B_2$は
四つの相異なる点からなるものの、
$\PGL_2$同型類集合。
\end{enumerate}
この二集合の間に自然な全単射がある。
対応は、$F_s$の二重被覆
\[
 \pi_F:E\longrightarrow\PP H^0(E,F_s)^\vee
\]
の四つの分岐点に対し、
$P$に対応する非零二捩れ平行移動が
誘導する二組の分割を取ることで与えられる。
逆方向には、四分岐点$B_1+B_2$の
二重被覆$\pi_F:E\to B$を取り、
その分割に対応する唯一の
非零$E[2]$の元を$P$とする。
このとき$F_s=\pi_F^*\OO_B(1)$である。

論文B
\cite[\texttt{thm:noncubic-intrinsic-presentation-equivalence}]{QuadricModuli2026}
の非cubic根対表示によれば、
従ってこの二捩れsquare壁の幾何的avatar同型類は、
\emph{四根が相異なるnoncubic O-marked H2境界}
の同型類とも全単射する。

最後の全単射は\emph{論文Bの根対分類を介した}
同型類の対応である。
この幾何辞書だけから
avatar側のAS$\mathbb Z$-代数の
$\qgr$とそのnoncubic O-marked heartの
同値、または対応するモジュライstackの射は
まだ結論しない。
\end{theorem}
\begin{proof}
次数二の線束$F_s$は楕円曲線から$\PP^1$への
二重被覆を与える。
その被覆対合を$\iota_F$と書くと、
その固定点集合$R_F$は四点からなり、
像$B_F$が四つの相異なる分岐点となる。
$E$は標識なしでも$\Pic^1(E)$のtorsorで、
$P\in\Pic^0(E)[2]\setminus0$は
$E$上の固定点のない位数二平行移動$t_P$を
定める。
次数二線束に対する平行移動の作用は
$t_P^*F_s\simeq F_sP^{\pm2}\simeq F_s$なので、
$t_P$は$\iota_F$と可換であり、
$R_F$を保つ。
その四点への作用は自由であり、
ちょうど二つの長さ二の軌道に分ける。
二つの軌道の$\pi_F$による像が
$B_F=B_1+B_2$の二分割となる。
これは線束の基底・楕円原点の選択に依存せず、
曲線と線束の同型に自然である。

逆に四つの相異なる分岐点$B_1+B_2$から、
$B$上の二重被覆は、
$\mathcal O_B(2)$とその四点を零点に持つ
切断を用いて、被覆同型を除き一意に構成される。
$F_s=\pi_F^*\OO_B(1)$は次数二である。
四つの分岐点は$E$上の被覆対合の固定点と
一対一に対応する。任意の一つを楕円原点に
選べば、その四点が$E[2]$となる。
$E[2]\simeq(\mathbb Z/2)^2$の各非零元の
平行移動は四点を二つの組に分割し、
三つの非零元と三つの二分割との間に
自然な全単射を与える。
従って指定の$\{B_1,B_2\}$から
一意の非零二捩れ$P$を回収する。
原点を変えても$\Pic^0(E)$における$P$の
同型類は変わらない。
二方向の構成が互いに逆であることも
この四点の分割から直ちに従う。

具体的に、適当な射影座標で
$B_1=\{0,\infty\}$、$B_2=\{1,\lambda\}$、
$\lambda\notin\{0,1,\infty\}$とすれば、
$E$はLegendre曲線
\[
 y^2=x(x-1)(x-\lambda)
\]
で与えられる。$B_1,B_2$をそれぞれ
集合として保ち、その二点を交換する
$\PP^1$の対合は$x\mapsto\lambda/x$である。
その固定点を持たないliftは
\[
 (x,y)\longmapsto
 \left(\frac{\lambda}{x},
             -\frac{\lambda y}{x^2}\right)
 \tag{Root-2}\label{eq:qa-avatar-root-pair-translation}
\]
であり、$E$上では非零二捩れ点
$(0,0)$によるtranslationである。
この式からも二分割と二捩れ点の
相互復元を直接確認できる。

最後に論文Bの非cubic被約境界は、
二つの二根集合の非順序対を
$\PGL_2$で割った群oidとして記述されており、
四根相異なる開ではまさに第二集合である。
従ってBの分類との合成が最後の
同型類全単射を与える。
これはBのO-marked対象への
逆構成による対応であり、
元のavatar AS代数の直接認識とは区別する。
\end{proof}

\subsection{二根対のquartic resolventと退化avatarの非可逆Picardパラメータ}
\label{subsec:qa-rootpair-resolvent-singular-avatar}

\begin{theorem}[noncubic根対の三重quartic resolventと三つの退化型]
\label{thm:qa-noncubic-rootpair-quartic-resolvent}
$k$を標数零の代数閉体、$W$を二次元ベクトル空間とし、
$\mathcal D_2=\PP(\operatorname{Sym}^2W^\vee)$、
$\mathcal D_4=\PP(\operatorname{Sym}^4W^\vee)$と置く。
二つの二次形式をその順序を忘れて掛ける射
\[
 \mu:\operatorname{Sym}^2(\mathcal D_2)\longrightarrow
             \mathcal D_4,\qquad
 \{[f_1],[f_2]\}\longmapsto[f_1f_2]
 \tag{Res-4}\label{eq:qa-rootpair-resolvent-map}
\]
は全射有限射であり、一般次数は$3$である。
根が相異なる四次形式$q$の上の三点は
四根の非順序二根対への三分割を表す。
アフィンmonic chartで
\[
 q(x)=x^4+a x^3+b x^2+c x+d
\]
とすると、四根$\alpha_1,\ldots,\alpha_4$の
二組分割から得る
$\alpha_i\alpha_j+\alpha_k\alpha_\ell$の
三つの値は
\[
 R_q(z)=z^3-bz^2+(ac-4d)z
             -(a^2d+c^2-4bd)
 \tag{Res-cub}\label{eq:qa-rootpair-resolvent-cubic}
\]
の三根に一致する。
この三重射は古典的なquartic resolventの
幾何的な表示である
\cite[Section 4.3]{Cremona1999}。

$\mathcal D_2^{\mathrm{sf}}$を二つの相異なる
根を持つ二次形式の開部分とする。
論文Bのnoncubic根対表示の
\[
 \mathcal R_{\mathrm{sf}}=
 \operatorname{Sym}^2(\mathcal D_2^{\mathrm{sf}})
\]
上では$q=f_1f_2$の幾何的根配置は
次の三型に限られ、それぞれ
\[
 C_q:\quad y^2=f_1f_2
 \quad\text{in }\operatorname{Tot}\OO_{\PP^1}(2)
 \tag{Res-curve}\label{eq:qa-rootpair-spectral-curve}
\]
の射影的二重被覆が次の形となる。
\[
 \begin{array}{c|c|c}
 \text{根配置}&q\text{の重複度}&C_q\\\hline
 B_1\cap B_2=\varnothing&(1,1,1,1)&\text{smooth genus }1\\
 |B_1\cap B_2|=1&(2,1,1)&\text{既約nodal genus }1\\
 B_1=B_2&(2,2)&\text{二つの}\PP^1\text{が二節点で交わる}
 \end{array}
 \tag{Res-strata}\label{eq:qa-rootpair-three-strata}
\]
三型とも$C_q$は連結なGorenstein曲線で
算術種数$1$である。
この構成は$\operatorname{PGL}(W)$と
根対交換の作用に自然であり、論文Bの
根対スタック
$\mathfrak R=[((\mathcal D_2^{\mathrm{sf}})^2)/
(\operatorname{PGL}(W)\times S_2)]$
からquartic二重被覆のスタックへの射を与える。
ただし全$\operatorname{Sym}^2(\mathcal D_2)$の
有限射と、平方因子を除いた
$\mathcal R_{\mathrm{sf}}$上の制限射は区別する。
\end{theorem}
\begin{proof}
四次形式の任意の因数分解は
その根の重複度付き四点集合から
二点ずつを選ぶ有限個の操作なので、
$\mu$の幾何fiberは有限である。
$\operatorname{Sym}^2(\mathcal D_2)$は射影的、
従って$\mu$はproperかつquasi-finiteで
有限である。代数閉体上ではすべての四次形式が
二次形式の積となるので全射。
四根が相異なれば分割は三つあり、
一般次数は三となる。
式~\eqref{eq:qa-rootpair-resolvent-cubic}は
四根の基本対称式を展開すれば得られる。
例えば三値を
$u_{12|34}=\alpha_1\alpha_2+\alpha_3\alpha_4$
などと置くと、
$\sum u=b$、
$\sum_{r<s}u_ru_s=ac-4d$、
$\prod u=a^2d+c^2-4bd$である。

$f_1,f_2$の根が各々相異なるので、
積の各根の重複度は高々二である。
共通根の個数が$0,1,2$の三場合が
表の各行に一致する。
四単根なら二重被覆はsmoothでRiemann--Hurwitzにより
genus一。一つの二重根の近くでは
$y^2=t^2u(t)$、$u(0)\ne0$なので、
étale局所で$y'^2=t^2$という普通節点になる。
残り二単根から二重被覆は既約で、
正規化はgenus零。
二つの二重根なら$f_1=f_2$（射影類として）で、
$y^2=f_1^2$は二つのsection
$y=\pm f_1$の和である。
各sectionは$\PP^1$で、$f_1$の二単根で横断する。
二重被覆の代数
$\OO_{\PP^1}\oplus\OO_{\PP^1}(-2)y$
から$\chi(\OO_{C_q})=1+(-1)=0$であり、
全型で$p_a(C_q)=1$。
局所的に平面曲線の節点またはsmoothなので
Gorensteinでもある。
群の作用については積と二重被覆の
構成が自然であることから従う。
\end{proof}

\begin{theorem}[根対からの相対torsion-free spectral avatar]
\label{thm:qa-noncubic-rootpair-torsionfree-spectral-module}
論文Bの二形式表示
\[
 (f_1,f_2)\in
 \bigl(\operatorname{Sym}^2W^\vee\bigr)_{\mathrm{nd}}^2
\]
上で、$B=\PP(W)$と置き、
$\widetilde p:\widetilde C\to B$、
$p:C\to B$を以下の有限射、
$r:\widetilde C\to C$をその間の自然な射とする。
\[
 \widetilde C=
 \operatorname{Spec}_{\PP(W)}
  \frac{\OO_{\PP(W)}[c,d]}
       {(c^2-f_1,\ d^2-f_2)},\qquad
 C=\operatorname{Spec}_{\PP(W)}
  \frac{\OO_{\PP(W)}[y]}{(y^2-f_1f_2)}
\]
を取る。ただし$c,d$は局所的に
$\OO_{\PP(W)}(-1)$の生成元、$y=cd$は
$\OO_{\PP(W)}(-2)$の生成元と解釈する。
$\widetilde C$から$C$への有限射は
$(c,d)\mapsto(-c,-d)$の商であり、
$\OO_C$-加群としての反不変部分と
その$B$への直像はそれぞれ
\[
 \mathcal M=(r_*\OO_{\widetilde C})^-,
 \qquad
 p_*\mathcal M
 \cong\OO_B(-1)c\oplus\OO_B(-1)d
 \tag{Res-M}\label{eq:qa-rootpair-spectral-module}
\]
であり、$\mathcal M$は
\[
 y c=f_1 d,\qquad y d=f_2 c
 \tag{Res-MF}\label{eq:qa-rootpair-spectral-matrix-factorization}
\]
を満たす。全二形式族と任意の基底変更にわたり、
$C$は有限平坦二重被覆、$\mathcal M$は
\emph{相対平坦なrank-one torsion-free sheaf}
（各幾何fiber上で次数零）となる。
この構成は二形式の基底変更・スカラー変更・
二形式交換に関して自然なので、
論文Bのnoncubic表示スタック
$\mathfrak N^{\mathrm{pres}}=[P/G]$から
「種数一の有限二重被覆とrank-one
torsion-free sheaf」の相対群oidへの射を与える。

二根対が交わらない滑らかなfiberでは、
$\widetilde C\to C$はétale $\mu_2$-torsor、
$\mathcal M$は非自明な
$\Pic^0(C)[2]$の線束である。
これは定理
~\ref{thm:qa-avatar-two-torsion-root-pair-bijection}
で定めた分割の2-torsionと一致する。
他方、根対が一根を共有するとき
$\mathcal M$はその節点で局所自由でなく、
二根対が一致するときは二つの節点で
局所自由でない。
従って\emph{全noncubic根対配置にわたって}
自然に存在するのは通常の$\Pic^0$線束ではなく、
この相対torsion-free spectral moduleである。
\end{theorem}
\begin{proof}
二つの二次平方根の代数は、
$\OO_B$上の基底$1,c,d,cd$を持つ
階数四の局所自由代数である。
$\widetilde p_*\OO_{\widetilde C}$の
同時符号変換の不変部分は$1,cd$、
反不変部分は$c,d$なので
式~\eqref{eq:qa-rootpair-spectral-module}を得る。
$y=cd$による積は
式~\eqref{eq:qa-rootpair-spectral-matrix-factorization}
で与えられ、形式的には
\[
 \begin{pmatrix}0&f_2\\f_1&0\end{pmatrix}^2
      =f_1f_2 I_2
\]
というmatrix factorizationである。
$p_*\mathcal M$は$B$上で
$\OO(-1)^{\oplus2}$だから、
二形式パラメータとその任意基底変更に対して平坦。
各幾何fiberでは二次拡大のgeneric点でrank一で、
局所的に二次平方根代数の直和因子であるから
torsion-freeである。
Euler特性は
\[
 \chi(C,\mathcal M)
 =\chi(\PP^1,\OO(-1)^{\oplus2})=0
 =\chi(C,\OO_C)
\]
なので次数零である。

$f_1,f_2$に共通零点がなければ、
$c,d$が同時に零となる点がないため
同時符号対合は自由で、
$\widetilde C\to C$は有限étaleの二重torsor。
このとき反不変部分$\mathcal M$は
invertibleかつ$\mathcal M^2\simeq\OO_C$。
$H^0(C,\mathcal M)=H^0(B,p_*\mathcal M)
 =H^0(\PP^1,\OO(-1)^{\oplus2})=0$
なので非自明である。
$\widetilde C$の三つの二次部分被覆の
分岐分割から、これは前定理の2-torsionに一致する。

共有根の局所では単元の平方根をétale局所に取り、
\[
 \OO_{C,o}^{\wedge}
  \simeq k[[x,y]]/(y^2-x^2),\qquad
 \mathcal M_o\cong
   \langle c,d\mid yc=xd,\ yd=xc\rangle .
\]
$\mathcal M_o/(x,y)\mathcal M_o$の次元は二なので、
rank一局所自由ではない。
一致根対では$C$は二つのsectionの和で、
その二つの交点の局所方程式も同じになり、
$\mathcal M$は各節点で非局所自由である。
二形式の交換は$c,d$の交換、
スカラー変更は平方根のスカラー変更として
作用し、構成は元の表示群の作用に整合する。
\end{proof}

\begin{corollary}[noncubic中央二重平方環の可換点曲線からの幾何射]
\label{cor:qa-noncubic-central-square-point-curve-spectral-avatar}
論文A・Bのnoncubic中央二重平方環
\[
 B_f=k[u,v]\langle c,d\rangle/
          (c^2-f_1(u,v),\,d^2-f_2(u,v))
\]
に対し、通常の可換化を施した環と
その射影曲線は
\[
 B_f^{\mathrm{ab}}=
 k[u,v,c,d]/(c^2-f_1,\,d^2-f_2),\qquad
 E_f=\operatorname{Proj}(B_f^{\mathrm{ab}})
 \subset\PP^3
\]
である。これは定理
~\ref{thm:qa-noncubic-rootpair-torsionfree-spectral-module}
の四重被覆$\widetilde C$と自然に同型であり、
\[
 \boxed{\quad
 E_f\xrightarrow{/\,\mu_2}C_{f_1f_2}
 \xrightarrow{\;2:1\;}\PP^1,
 \qquad
 \mathcal M=(r_*\OO_{E_f})^- .
 \quad}
 \tag{NC-point-avatar}\label{eq:qa-noncubic-pointcurve-spectral-factorization}
\]
$\mu_2$は$(c,d)\mapsto(-c,-d)$として作用する。
この対合の固定点は$f_1,f_2$の
共通根の上に限られるので、二根対が
非交差なら第一射はétale、
一根共有または一致根対では
固定点に対応する節点でétaleでなくなる。
この場合、有限平坦な二重torsorと
みなしてはならない。

この射はnoncubic線形代数$B_f$から
rank-one torsion-free spectral avatarへの
\emph{具体的な幾何構成}である。
非可換な$B_f$の$\qgr$と
$\OO$-marked avatar側のheartが同値である
と主張するものではない。
\end{corollary}
\begin{proof}
$u,v$は中央にあるので、可換化において
追加される関係は$cd=dc$だけである。
そのProJの$\PP^1=\operatorname{Proj}k[u,v]$上の
相対座標環は$1,c,d,cd$を基底とする
二重平方根代数となる。
これは前定理の$\widetilde C$そのものである。
同時符号の不変部分は$1,cd$、
反不変部分は$c,d$なので、
商と$\mathcal M$の表示も一致する。
非自明な同時符号の固定点は、
射影座標$(u,v)$が非零である限り
$c=d=0$を満たす点に限られ、
これは$f_1=f_2=0$と同値である。
$(u,v)=(0,0)$では二つの平方関係から
$c=d=0$となり射影点は存在しない。
従って固定点の条件が従う。
\end{proof}

\begin{proposition}[共有根退化での可逆Picard延長の局所障害]
\label{prop:qa-rootpair-nodal-picard-nonextension}
$R=k[[s]]$上の二根対の退化
\[
 f_1=x(x-1),\qquad f_2=(x-s)(x-2)
\]
を取り、$C_R:y^2=f_1f_2$とする。
その一般fiberは滑らかな楕円曲線であり、
特殊fiberは既約な一節点曲線となる。
一般fiberの自然な非零2-torsion線束
$\mathcal M_K$は、
\emph{$C_R$全体の可逆線束には延長できない}。
従って二根対から
$C:y^2=f_1f_2$と
\emph{可逆な}$P\in\Pic^0(C)$を
割り当てる滑らかなavatar辞書は、
この族に沿ってそのまま相対射としては延長できない。

この障害は、$C_R$の節点における
表面$A_1$特異点と、
そこでの非自明なrank-one reflexive module
によって検出される。
$C_R$以外の曲線へのavatar再埋込み、
$\Sigma_2$へのelementary transform、
またはcompactified Picard型データ
\cite{AltmanKleiman1980}による
延長の可能性は排除しない。
ただし同じ一径数基底の上で全空間の$A_1$特異点を
最小解消するだけでは不十分である。
命題~\ref{prop:qa-nodal-resolution-odd-exceptional-degree}の
反不変moduleの可逆化は例外曲線上の次数が奇数であり、
任意の垂直Cartier因子でtwistしても
相対的な多重次数零線束にはならない。
第~\ref{sec:qa-spectral-boundary-models}節の二traitでは，
任意の有限DVR基底変更後の最小正則模型でも二乗自明な線束への延長は存在しない
（定理~\ref{thm:qa-boundary-input-two-torsion-line-obstruction}）．
線束自体への延長と，非可逆TF層・二次積を保持する標的は区別する．
\end{proposition}
\begin{proof}
$s=0$では$f_1f_2=x^2(x-1)(x-2)$であり、
特殊曲線は既約nodal genus一である。
共通根$x=s=0$の完成局所環では、
単元の平方根で座標を正規化すると
\[
 A=k[[x,s,y]]/(y^2-x(x-s)).
\]
二形式を別々に平方根化する被覆の完成環は
$B=k[[c,d]]$で、
\[
 x=c^2,\quad x-s=d^2,\quad y=cd,\quad
 A\simeq B^{\{\pm1\}}
\]
となる。$B$上の同時符号の反不変部分
$M=Ac+Ad$は、商$A$上の
rank-one maximal Cohen--Macaulay、従って
rank-one reflexive moduleである。
その最小生成元は$c,d$の二つなので非自由であり、
$A$は表面$A_1$特異点を持つ。
これは前定理の$\mathcal M$の完成stalkである。

相対曲線の全空間$C_R$はnormalであり、
その唯一の特異点で上記$A_1$局所型を持つ。
$\mathcal M$はnormal曲面上の
rank-one reflexive sheafで、
一般fiberでは$\mathcal M_K\in\Pic^0(C_K)[2]$。
もし可逆線束$L$が全空間上に存在して
$L_K\simeq\mathcal M_K$となれば、
normal曲面上のrank-one reflexive sheafの
Weil因子記述から、$L$と$\mathcal M$の差は
特殊fiberに支持されるWeil因子となる。
この特殊fiberは\emph{既約}かつ
パラメータ$s$で切られたCartier因子なので、
その整数倍は主因子である。
従って$L\simeq\mathcal M$となり、
節点での非局所自由性に矛盾する。
最後の留保は$C_R$を固定したときの
可逆延長障害と、別のavatarモデルの存在を
混同しないためである。
\end{proof}

\begin{theorem}[全noncubic根対と二捩れsquare avatarの体上同型類全単射]
\label{thm:qa-noncubic-rootpair-all-geometric-square-avatar-bijection}
$k$を標数零の代数閉体とする。
論文Bのnoncubic根対の全幾何点、すなわち
$\PP^1$上の\emph{各々は被約}な二点因子の
非順序対$\{B_1,B_2\}$の
$\PGL_2$-同型類集合を考える。
これは、次の条件を持つsquare avatarの
\emph{同型類集合}と自然に全単射する。
\[
 \begin{gathered}
 (Y=\PP^1\times\PP^1;\
  D\in|\OO_Y(2,2)|;\ F_s,G_s,P),\\
 P=G_sF_s^{-1}\ne\OO_D,\qquad
 P^{\otimes2}\simeq\OO_D,\\
 D\text{ はsmooth genus }1,\quad
 \text{既約一節点、または二}\PP^1
 \text{が二節点で交わる曲線}.
 \end{gathered}
 \tag{Sq-all-nc}\label{eq:qa-all-noncubic-square-avatar-bijection}
\]
ここで$F_s=\OO_D(1,0)$、
$G_s=\OO_D(0,1)$であり、
二rulingの順序は保持する。
この全単射は、論文Bの
\cite[\texttt{thm:noncubic-intrinsic-presentation-equivalence}]{QuadricModuli2026}
を介して、\emph{全noncubic O-marked H2幾何同型類}の
集合論的なavatar分類を与える。

構成は次の通りである。
まず$q=f_1f_2$から
式~\eqref{eq:qa-rootpair-spectral-curve}の
二重被覆$p:C_q\to\PP^1$と
$F_s=p^*\OO_{\PP^1}(1)$を取る。
\begin{enumerate}
\item 四根相異なる場合は、
定理~\ref{thm:qa-avatar-two-torsion-root-pair-bijection}
の分割に対応する非零
$P\in\Pic^0(C_q)[2]$を取る。
\item 一根共有または二根対一致の場合は
$\Pic^0(C_q)\simeq\mathbb G_m$なので、
その\emph{唯一の}非自明二捩れ点$P=-1$
を取る。
\end{enumerate}
いずれの場合も$G_s=F_sP$は次数二で
大域生成され、二つの完全線形系の積
\[
 C_q\longrightarrow
 \PP H^0(F_s)^\vee\times\PP H^0(G_s)^\vee
\]
は$(2,2)$反標準Cartier因子への閉埋込みとなる。
この逆操作は、square avatarの
第一rulingの二重被覆$D\to\PP^1$の
分岐四次因子と$P$から、
四根相異なる場合には三分割の一つを、
一根共有・一致根対では唯一の
被約二根対分割を復元することである。

\emph{重要：} これは\emph{体上の同型類の全単射}である。
上の二場合で特別に選んだ
$P=-1\in\mathbb G_m$は、
必ずしも滑らかな族の2-torsion線束の極限でない。
実際、命題
~\ref{prop:qa-rootpair-nodal-picard-nonextension}
の一根共有族では延長できない。
従ってこの全単射を、そのまま
通常のPicard avatarスタックへの
相対モジュライ射、または元のavatar
AS代数の$\qgr$のpointed同値と
読み替えてはならない。
\end{theorem}
\begin{proof}
smoothな場合は定理
~\ref{thm:qa-avatar-two-torsion-root-pair-bijection}
による。以下、特異な二場合を調べる。

\emph{既約一節点の場合。}
$C_q$の正規化は$\nu:\PP^1\to C_q$で、
節点の二逆像を$a,b$とする。
$\Pic^0(C_q)$の線束は、
正規化上の自明線束の
$a,b$での貼合せスカラー
$\lambda\in k^\times$で分類される。
従って$\Pic^0(C_q)\simeq\mathbb G_m$で、
非自明二捩れは貼合せ比$-1$だけである。

$F_s=p^*\OO(1)$は正規化上次数二で、
節点での貼合せ比を$1$に規格化できる。
$G_s=F_sP$は正規化上の$\OO(2)$で、
貼合せ比は$-1$となる。
$H^0(C_q,G_s)$は$H^0(\PP^1,\OO(2))$の
$a,b$における値が符号反対になる
二次元部分空間である。
これはbasepoint-freeである。
実際、$a,b$の二評価汎関数の
非自明な線形結合は、
二次Veroneseの双対円錐上で
単一点の評価汎関数には一致しない。
（異なる二点を結ぶ割線が二次曲線と
交わるのはその両端のみである。）
従って$|G_s|$も有限な次数二射を与える。

$F_s\not\simeq G_s$なので、
この二射の積は総次数二ではあり得ず、
$C_q$上の一般点で単射である。
像の二rulingへの次数は各二だから
$(2,2)$因子で、算術種数は一。
元の$C_q$も算術種数一なので、
有限双有理射の構造層の余核はEuler特性零、
従って零であり、積射は閉埋込みとなる。

\emph{二根対一致の場合。}
$C_q=C_+\cup C_-$は二つの$\PP^1$が
相異なる二節点で交わる曲線である。
正規化上の各成分の次数零線束は自明であり、
二つの節点での貼合せ比
$(\lambda_1,\lambda_2)$を
成分ごとの定数倍で同一視すると、
$\lambda_1/\lambda_2\in\mathbb G_m$
だけが残る。
よって$\Pic^0(C_q)\simeq\mathbb G_m$、
非自明二捩れはこの比が$-1$のものに限る。
$F_s$は各成分で$\OO_{\PP^1}(1)$であり、
$G_s=F_sP$も各成分で次数一、
二節点での値の二条件から
$H^0(C_q,G_s)$は二次元で大域生成される。
二つの成分の像は、
第一$\PP^1$から第二$\PP^1$への
二つの$(1,1)$グラフとなる。
貼合せ比$-1$ゆえ二グラフは異なり、
ちょうど二つの節点で横断する。
従って像は元の$C_q$と同型な
$(2,2)$Cartier因子である。

逆向きは平方freeな二つの二次因子の
分割を数えることで得られる。
四単根なら三通りで非零$E[2]$の三点に対応する。
$(2,1,1)$なら
$(a,b)\mid(a,c)$の一通り、
$(2,2)$なら
$(a,b)\mid(a,b)$の一通りしか
二因子ともsquarefreeにする分割はない。
同じ各曲線の$\Pic^0[2]\setminus0$の
点数が$3,1,1$なので、逆は一意である。
元の第二rulingは
$G_s=F_sP$で復元されるから、
曲線・線束の同型はavatarの同型となる。
最後にBのnoncubic根対分類と合成すれば
O-marked同型類への全単射が得られる。
相対射に関する留保は直前命題の
非延長例により必要である。
\end{proof}

\begin{theorem}[全noncubic O-marked同型類の$\Sigma_2$反標準avatar分類]
\label{thm:qa-all-noncubic-sigma2-anticanonical-avatar-bijection}
$k$を標数零の代数閉体とする。
$\Sigma_2=\PP_{\PP^1}(\OO\oplus\OO(2))$の
負断面を$S=h-2f$とし、
$D\in|2h|=|-K_{\Sigma_2}|$を
$D\cap S=\varnothing$を満たす
相対反標準Cartier因子とする。
このとき、次の二集合は自然に全単射する。
\[
 \boxed{
 \begin{aligned}
 &\{\text{全noncubic O-marked quadricの
             幾何exact同型類}\}\\
 &\quad\longleftrightarrow
 \left\{
 \begin{array}{l}
 (D\subset\Sigma_2,P),\quad
 D\cap S=\varnothing,\\
 0\ne P\in\Pic^0(D)[2]
 \end{array}
 \right\}
/\operatorname{Aut}(\Sigma_2;h,f).
 \end{aligned}}
 \tag{NC-Sigma2}\label{eq:qa-all-noncubic-sigma2-avatar}
\]
右辺で$D$はあらかじめ滑らか、既約、被約と
仮定しない。しかし$P$の非自明二捩れ性から、
実際には$D$はsmooth genus一、
既約一節点、または二つの$\PP^1$が二節点で
交わる型のいずれかになる。
この三型は論文Bの二根対の
非交差／一根共有／一致に正確に対応する。

具体的な同型類の対応は
定理~\ref{thm:qa-noncubic-rootpair-all-geometric-square-avatar-bijection}
のsquare avatar
$(D_s;F_s,G_s,P)$に、
系~\ref{cor:qa-square-chart}の
幾何的変換$\mathsf T^{-1}$を適用したものである。
$\Sigma_2$側の制限線束を
$F=\OO_{\Sigma_2}(f)|_D$、
$H=\OO_{\Sigma_2}(h)|_D$とすると、
\[
 H\simeq F^2,\qquad
 F_s=F,\quad G_s=FP,\quad
 H_s=F_sG_s=F^2P
 \tag{NC-Sigma2-dict}\label{eq:qa-noncubic-sigma2-square-dictionary}
\]
が全対応の明示公式となる。

これは\emph{体上の集合論的全単射}であり、
$\Sigma_2$ avatarから元のAS代数の
pointed heartを直接回復する定理でも、
Picard line bundleを使う相対スタック同値でもない。
smooth域でこのsquare換算を指定AS10 mutationと一番頂点の
pointingへ接続した対象は、
定理~\ref{thm:qa-spectral-avatar-pointed-involution}により
元の対象の二同種双対に対応する。
元とpointed同値となるのは$\chi=0,2$だけであり、
この対応先は二つのnoncubic境界でH2族にも延長できない
（系~\ref{cor:qa-spectral-involution-no-h2-extension}）。
命題~\ref{prop:qa-rootpair-nodal-picard-nonextension}
の共有根退化では、固定されたspectral曲線族上の
可逆$P$への相対延長が不可能なことに注意する。
\end{theorem}
\begin{proof}
$D\cap S=\varnothing$のとき、
$S=h-2f$の定義から
$\OO_D(S)\simeq\OO_D$となり、
$H\simeq F^2$が従う。
また$D$はどの$\PP^1$-fiberも成分として含まない。
実際、fiberは$S$と交わるので、
そのような成分は$D\cap S=\varnothing$に反する。
従って$p:D\to\PP^1$は有限平坦次数二である。
そのtrace-zero直和分解を取ると
\[
 p_*\OO_D\simeq\OO_{\PP^1}\oplus
                \OO_{\PP^1}(-2)
\]
となる。後項の次数は
$\chi(\OO_D)=0$と$p_a(D)=1$から決まる。
よって$D$は四次式$q$を用いた
$y^2=q$型の二重被覆になる。
$q$が四単根、一つの二重根、二つの二重根なら
前定理の三型になる。
それ以外の四次根配置は
$(3,1)$のcuspidal型、
$(4)$の二枝接触型、
または$q=0$の非被約double sectionである。
これらの$\Pic^0$は標数零で
加法型$\mathbb G_a$となり、
非自明二捩れ点を持たない。
この主張は正規化と特異点での単元の貼合せ、
非被約double sectionでは
$1+\OO_{\PP^1}(-2)$の加法的拡大から得られる。
従って$P\ne\OO_D,\ P^2=\OO_D$は
三つのsmooth/nodal/I2型を強制する。

幾何的変換については
系~\ref{cor:qa-square-chart}が、
$\Sigma_2$の負断面を避けるavatarを
square chartへ送る操作$\mathsf T$を
任意基底変更と可換な同値として構成している。
そこで$H=F^2$を代入すると
square側の線束は
\[
 F_s=F,\qquad G_s=HPF^{-1}=FP
\]
となり、$P=G_sF_s^{-1}$かつ$P^2=\OO_D$である。
逆にこの二捩れsquare wallのデータに
$\mathsf T^{-1}$を適用すれば
$H=F_sG_sP^{-1}=F_s^2$となり、
系~\ref{cor:qa-square-chart}と
定理~\ref{thm:qa-open-charts}により
負断面を避ける$\Sigma_2$に入る。
両操作は互いに逆であり、
曲面・ruling・因子・Picard類の同型を保つ。
これと前定理の全noncubic二根対の
幾何同型類全単射、論文Bの
\cite[\texttt{thm:noncubic-intrinsic-presentation-equivalence}]{QuadricModuli2026}
を合成して結論を得る。
相対射に関する留保は直前の
$A_1$局所非延長の結果による。
\end{proof}

\begin{corollary}[noncubic良moduliの同じ点を区別する二つの$\Sigma_2$ avatar]
\label{cor:qa-sigma2-avatar-separates-special-noncubic-classes}
論文Bのnoncubic境界の良moduli
$\chi:\mathfrak N\to\mathbb A^1$において、
一根共有と二根対一致はともに$\chi=1$へ写るが、
これらは異なるpointed exact同型類である。
定理~\ref{thm:qa-all-noncubic-sigma2-anticanonical-avatar-bijection}
の$\Sigma_2$ avatarによれば、
前者は\emph{既約nodal}反標準因子$D$と
唯一の非零二捩れ$P$、
後者は\emph{二$\PP^1$が二節点で交わる}
反標準因子$D$と唯一の非零二捩れ$P$に対応する。
従って曲線$D$の既約成分構成が、
良moduliでは同一視される二つの
O-marked同型類を区別する。
\end{corollary}
\begin{proof}
良moduliの特殊fiberの二幾何同型類と
$\chi=1$への写像は論文B
\cite[\texttt{prop:noncubic-good-moduli-fibers}]{QuadricModuli2026}
による。
根対の共通根の個数に応じた
二重被覆$C_q$の特異型の分類は
定理~\ref{thm:qa-noncubic-rootpair-quartic-resolvent}
にあり、$\Sigma_2$での再埋込みは
前定理で証明した。
既約性・成分数が異なるので
二avatarは同型になり得ない。
\end{proof}

\begin{remark}[節点型avatarの二つの明示$(2,2)$模型]
\label{rem:qa-singular-square-avatar-explicit-models}
定理~\ref{thm:qa-noncubic-rootpair-all-geometric-square-avatar-bijection}
の特異な二型は、いずれも
$\PP^1_X\times\PP^1_Z$上で明示できる。
既約一節点型では、正規化$\PP^1_t$の
$t=1,-1$を貼り合わせ、非自明二捩れの
貼合せ比を$-1$とする。
$F_s$の二切断は$1,t^2$、
$G_s$の二切断は$t,1-t^2$であり、
\[
 X=t^2,\qquad Z=\frac{1-t^2}{t}
 \quad\Longrightarrow\quad
 XZ^2=(1-X)^2 .
 \tag{Sq-node}\label{eq:qa-nodal-square-avatar-model}
\]
これは$(X,Z)=(1,0)$に普通節点を持つ
既約$(2,2)$曲線である。
式は各$\PP^1$で斉次化するものとする。

一致根対型では、$C_+\cup C_-$を二つの
$\PP^1$の和として、その二交点を
$X=0,\infty$に取る。$F_s$は両成分で
$X$座標、非自明二捩れ$P$による
$G_s$は第一成分で$Z=X$、
第二成分で$Z=-X$となる。
従って
\[
 (Z-X)(Z+X)=0
 \tag{Sq-I2}\label{eq:qa-i2-square-avatar-model}
\]
は二つの$(1,1)$曲線が
$(0,0),(\infty,\infty)$で横断する
反標準$(2,2)$因子を与える。
この二模型は、特異な場合の
$P=-1\in\Pic^0(D)[2]$が本当に
square avatarの二rulingを作ることを
直接に示している。
\end{remark}

\begin{proposition}[Picard零切断の全体を古典pointとして追加できない理由]
\label{prop:qa-avatar-zero-picard-collapse}
$Y=\PP^1\times\PP^1$、
$D\in|\OO_Y(2,2)|$を任意の反標準Cartier因子とし、
$P=\OO_D$と置く。
このとき曲線上の四線束
\[
 (\OO_D,\ F_s,\ G_s,\ F_sG_s)
\]
の有向End代数（各積を含む）は、
$D$の選択に依存せず、
$Y$上の可換quadricの標準四項collection
\[
 (\OO_Y,\OO_Y(1,0),\OO_Y(0,1),\OO_Y(1,1))
\]
の有向End代数と同型である。
従って$\Pic^0$の零切断をそのまま
曲線型avatarに追加すると、
異なる$D$が同一の可換foundation代数を与え、
O-marked分類への単射性が失われる。
古典的な可換quadricは、曲線型とは別に
一点として追加する必要がある。
\end{proposition}
\begin{proof}
相異なる四項の間のHomに現れる
$\OO_Y(a,b)$は、
$a,b\in\{-1,0,1\}$の範囲にある。
各々についてCartier制限列
\[
 0\to\OO_Y(a-2,b-2)\to
 \OO_Y(a,b)\to\OO_D(a,b)\to0
\]
を考えると、$a-2\le-1$、$b-2\le-1$であり、
Künnethの公式から
$H^0(Y,\OO(a-2,b-2))=
H^1(Y,\OO(a-2,b-2))=0$である。
従って四項間の全Homの制限は同型となる。
構成は線束の合成に関して自然なので、
有向End代数の積も保存する。
右辺は$D$に依存しない。
\end{proof}

\begin{remark}[二つの $\Sigma _2$ への移行]
\label{rem:qa-rotated-square}
square helix の foundation を一つ進めた rotated square の偏極は
\[
 F_r=F_s,\qquad H_r=F_sG_sP
\]
である（本稿の幾何的証明は定理~\ref{thm:qa-core-avatar-square-rotated}；
原草稿\cite[命題11.1・定理15.1]{AvatarVSeven2026}とも整合する）。
一方、\eqref{eq:qa-square-dictionary}の偏極は $H=F_sG_sP^{-1}$ である。
前者の $\Sigma _2$ 条件は $P=F_sG_s^{-1}$、後者の条件は $P=G_sF_s^{-1}$ となる。
これらは一般には別の壁である。
foundation の移動や mutation とともに選ばれる構造対象も変わるため、
固定された構造対象の内在的一点 Hilbert 関手を同じものとして扱わない。
\end{remark}

\subsection{古典点と最深境界の扱い}

本節の stack は一貫して $P$ が各 fiber で非自明な曲線型を扱った。
体上の分類では、$P=\OO_D$ とした有限 directed algebra は $D$ の選択に依存せず、
\[
 \mathrm{Cl}_0:\operatorname{coh}(\PP^1\times\PP^1),\qquad
 \mathrm{Cl}_2:\operatorname{coh}(\PP(1,1,2))
\]
という二つの古典型に対応する
\cite[注意11.3・定理15.3]{Q12Avatar2026}。
後者の $\PP(1,1,2)$ は重み付き射影 stack である。
これは同型類の説明であり、相対 Picard の零切断を二点へ自動的に縮約する
stack の構成ではない。零切断では $D$ とその自己同型が残るため、
相対的に古典点を加えるには別途の収縮または同一視を定義する必要がある。

最深の十二本矢層についても、三頂点の有限安定表現のモジュライが $\Sigma _2$ になることと、
標準 heart における内在的一点 Hilbert 曲面が $\Sigma _2$ になることは異なる。
この区別のための内在的な検査として，
例~\ref{ex:qa-weighted-not-h2}では重み付き古典模型の全標準頂点で，
例~\ref{ex:qa-hirzebruch-not-h2}では $\coh\Sigma_2$ 自身の全可逆層で，
反Serre軌道の生成条件が破れることを示した。
古典型の負断面上には、一つの foundation では安定であっても、
次の foundation では最後の頂点から生成されなくなる表現が存在する
\cite[命題14.1]{Q12Avatar2026}。
したがって本節の開部分 stack だけから、全 $\mathcal A_{12}^{\times}$ の
内在的 quadric モジュライへの射、あるいはその射の定義域が開であることを結論しない。
そこでは構造対象、標準 heart、豊富性、Serre 平方根および相対 AS 持上げを
別に検証しなければならない。

% !TEX root = ../nc-quadric-surface-avatar.tex
% 2026-11-23（canonical square完成の相対tails認識を含む）
\section{内在的なモジュライへの忘却}
\label{sec:qa-intrinsic-forgetting}

有限 foundation、周期的 AS 代数、avatar、および構造対象だけを指定した
Abel 圏は、保持する構造も同型の意味も異なる。
本節では、前節までの開部分と内在的モジュライとの関係を、
既に構成した射と、相対的な認識を要する射とに分けて述べる。

\subsection{標的となる強い内在的クラス}

本稿の \(\HH\)、\(\HHstack\) は
\cite[\texttt{def:h2}]{QuadricIntrinsic2026} と
\cite[\texttt{def:relative-h2}]{QuadricModuli2026} の意味で用いる。
特に、任意の三次元 AS 正則 \(\mathbb Z\)-代数の tails 圏を
無条件に \(\HH\) と呼ぶわけではない。

\begin{definition}[体上の強い内在的クラス]
\label{def:qa-h2}
\(\HH\) の対象は \((\mathcal C,\OO)\) であり、次を満たす。
\begin{enumerate}
\item \(\mathcal C\) は \(\kk\)-線形 noetherian Abel 圏で、
\(D^b(\mathcal C)\) は Hom 有限で Serre 関手 \(S\) を持つ。
\item \(\OO\) は rank-one torsion-free exceptional 対象である。
ここで $\mathcal C_0$ は有限長対象の部分圏、$\mathcal C_1$ は
$\mathcal C/\mathcal C_0$ で有限長となる対象の逆像とし、
$q:\mathcal C\to\mathcal C/\mathcal C_1$ とする。
条件は $q\OO$ が非零単純かつ $\Hom(\mathcal C_1,\OO)=0$ である。
また exceptional とは $\End(\OO)=\kk$、$\Ext^{i}(\OO,\OO)=0$ ($i>0$) をいう。
\item 完全自己同値 \(s\) が存在し、
\(s^2\simeq S^{-1}[2]\) かつ \((\OO,s)\) は Artin--Zhang ample である。
\item 同じ \(s\) に対して、全軌道環
\[
 B(\mathcal C,\OO,s)
   =\bigoplus_{n\ge0}\Hom(\OO,s^n\OO)
\]
が \(A/(g)\) と同型である。
ここで \(A\) は左右 noetherian、二次 Koszul AS 正則、
型 \((1,4,6,4,1)\) の環、\(g\in A_2\) は左右正則正規元である。
\end{enumerate}
\(s\)、根同型、\(A\)、\(g\)、商表示は指定構造ではない。
射は \(\OO\) の同型を伴う完全同値であり、
その同型を保つ自然同型によって一射を同一視する。
\end{definition}

ここで型 \((1,4,6,4,1)\) は、自明加群の最小解消の階数列が
この列で、global dimension と Gorenstein parameter がともに四であることを意味する。
\(g\) の中心性は要求しない。
また、ある商圏との同値だけでなく、指定した \(\OO\) と同じ \(s\) の
\emph{全軌道環} に対する表示を要求する。

相対版 \(\HHstack\) は、標数零の代数閉体上の局所有限型スキームの
fppf site 上で定義する。
圏の基底変更は Ind 圏における加群対象の圏によるものであり、
Hom を tensor するだけの基底変更ではない。
対象と完全列の fppf descent、有限 Hom と flat base change、
平坦な \(\OO\) を要求する。
局所 canonical compact 導来圏は相対 smooth proper で、
相対 Serre 双対性と導来基底変更に適合する。
全幾何ファイバーは定義~\ref{def:qa-h2} を満たす。
さらに、\(s\)、AZ 条件、全軌道環の AS4 被覆は
\emph{非被約基底も含めて fppf 局所に存在する}ことを要求する。
根や被覆を重なり上で同一視することは要求しない。
これらは標識ではなく存在条件だからである。

\begin{remark}[相対存在は幾何点条件より強い]
\label{rem:qa-relative-root}
全幾何ファイバーが \(\HH\) に入っても、
相対平方根、AZ ampleness、相対 AS4 被覆が存在することは別問題である。
例外対象の無障害性や有限代数の \(\operatorname{HH}^3\) の消滅だけでは、
これらの条件は得られない。
\cite[\texttt{prop:chosen-root-does-not-lift}]{QuadricModuli2026} は、可換 quadric 上の
ある適合根が、特定の一次 Abel 変形へ持ち上がらない例を与える。
これは別の根の存在まで否定する例ではなく、
\(\HHstack\) の非代数性や許容部分の非開性を証明するものでもない。
\end{remark}

\subsection{canonical square 完成で確立した定義域}

定理~\ref{thm:qa-relative-square-completion-tails} と
命題~\ref{prop:qa-relative-square-unfolding} により，実際の曲線 avatar
から構成した有限 square 代数の canonical 完成と全整数展開には，
標準対象を保つ相対 tails 認識がある．これは幾何点の同値だけではなく，
本稿の site 上の module-object 基底変更と coherent に整合する完全同値である．
左右 Noether 性と係数束付きの実際の AS 解消も証明済みである．
補題~\ref{lem:qa-classified-fiber-finite-square-recognition} は，
古典型を含む分類済み square fiber を持つ有限局所自由族まで扱う．
従ってこの canonical 部分の相対 heart 回復を未解決のままにしない．
一方，任意の別途与えた AS/QP 族との全整数同一視には
注意~\ref{rem:qa-square-unfolding-vertex-lines} の頂点線束・周期比較が必要であり，
以下の一般 presentation stack 全体の認識と区別する．

さらに，定理~\ref{thm:qa-square-normal-foundation-stack-equivalence} は
有限 square と normal foundation の全スタックを同定し，
系~\ref{cor:qa-h2-square-foundation-open} は既存の内在的開をこの表示へ移す．
全 $\HHstack$ については
定理~\ref{thm:qa-universal-h2-rank81-completion} が，ruling の選択を要しない
階数八十一の通常代数と pointed tails を与える．
ここまでの有限表示・canonical 完成の認識は確立済みである．
定理~\ref{thm:qa-square-explicit-h2-open}と
系~\ref{cor:qa-square-h2-twenty-polynomials}は，この全$H_2$開を
三十二次元線形表示上の四組・計二十式以内で記述する．
任意の別途与えた AS/QP presentation 群oidとの比較は後続課題である．

\subsection{忘却射の正確な定義域}

\begin{definition}[内在的条件に適合する族]
\label{def:qa-compatible-chart}
\(\mathfrak U\) を、周期的代数または avatar から
基底変更と同型に関して自然に \((\mathcal C,\OO)\) を構成する
fppf パラメータスタックとする。
\(\mathfrak U_{\HH}\) は、その構成した族が
上記の \(\HHstack\) の相対条件を満たす対象の全部分スタックとする。
この定義において、\(\mathfrak U_{\HH}\) が代数的または開であることは
仮定も結論もしない。
\end{definition}

\begin{proposition}[適合族からの忘却]
\label{prop:qa-compatible-forgetting}
定義~\ref{def:qa-compatible-chart} の構成は射
\[
 \pi_{\mathfrak U}:\mathfrak U_{\HH}\longrightarrow\HHstack,
 \qquad u\longmapsto(\mathcal C_u,\OO_u)
\]
を与える。
二つの chart 間の変換が、基底変更と両立する pointed exact 同値を
与えるなら、二つの忘却射はその変換に沿って自然に二同型となる。
\end{proposition}
\begin{proof}
\(\mathfrak U_{\HH}\) の条件は fppf 局所的で、
pointed exact 同値と基底変更に安定である。
従って同型で飽和した全部分スタックを定める。
構成した圏の系、構造対象、その比較同型を保持して
残る標識を忘れれば、定義された標的への functor を得る。
二つの構成を結ぶ pointed exact 同値が任意の基底変更と整合する場合には、
それらの同値がそのまま二つの functor の二同型を与える。
\end{proof}

この命題は忘却の意味を明示するものであり、
未知の条件を有限個の方程式で認識したという定理ではない。
特に \(\mathfrak F_{\le12}\) は有限代数だけを分類するので、
そこから直接 \(\mathcal C\) を取ることはできない。
\(\operatorname{Perf}(A)\) の構成だけでは、元の曲面の heart は選ばれていない。
有限代数の標準的な加群 heart と、tilting によって移された曲面の heart を
同一視してはならない。

「十二本矢の適当な開部分から \(\HHstack\) への自然な射」という主張を
有限な表示で実現するためには、少なくとも次を示す必要がある。
\begin{enumerate}
\item 使用する periodic algebra、grading、foundation と heart を
相対的に構成し、全次数の平坦性、noetherian 性、基底変更を確認する。
\item 指定した \(\OO\) の相対例外性・rank 条件、相対 Serre 公式、
適合する完全平方根と AZ ampleness を確認する。
\item 同じ根の全軌道環が、相対 AS4 環の正則正規二次切断となることを確認する。
\item これらの条件を持つ部分関手が、考えている表示スタックの
不変開部分で表されることを証明し、表示変更の全ての射を降下させる。
\end{enumerate}
根と全軌道環が既に与えられた場合、固定軌道環への AS4 被覆を選ぶ空間は
\cite[\texttt{thm:relative-finite-as-cover}]{QuadricModuli2026} により、
\(\operatorname{Gr}(6,K)\) 内の有限表示局所閉 incidence として記述できる。
しかし、その被覆を忘れた fppf 像の開性・代数性や、根そのものの存在は
この有限判定の結論ではない。

\subsection{十二本矢の全体を送ることへの具体的な障害}

\begin{lemma}[反 Serre 軌道の生成は必要である]
\label{lem:qa-antiserre-necessary}
\((\mathcal C,\OO)\in\HH\) ならば、
\(t=S^{-1}[2]\) に対する \((\OO,t)\) も AZ ample である。
特に、非零対象 \(M\) に対して
\(\Hom(t^{-n}\OO,M)=0\) が全ての \(n\ge0\) で成立することはない。
\end{lemma}
\begin{proof}
定義の根 \(s\) を選び、\(E=\OO\) と置く。
\(B=A/(g)\) は次数一で生成される。
AZ の生成条件を \(sE\) に適用し、得られた全射に
\(s^{-1}\) を作用させると、\(E\) は
\(s^{-m}E\)（\(m>0\)）の有限直和の商になる。
次数一生成により各成分 \(s^{-m}E\to E\) は、
最後の段が \(s^{-1}E\to E\) である合成の線形結合である。
従って評価 \(\Hom(s^{-1}E,E)\otimes s^{-1}E\to E\) は全射となる。
Hom 有限性により、この評価の源は有限直和である。
これは \cite[\texttt{lem:ample-square-from-generation}]{QuadricIntrinsic2026}
の評価全射である。

任意の対象への AZ 生成全射で奇数指数 \(m\) の項があれば、
この評価全射を \(s^{-m}\) で移して \(s^{-(m+1)}E\) の
有限直和に置き換える。従って偶数指数の項だけで生成できる。
指数に任意の下限を課す場合は、先に対象を正に shift して
生成全射を作り、負の shift で戻す。
全射に対する Hom の最終的な全射性は
元の AZ 条件を偶数部分列に制限すれば従う。
\(s^2\simeq t\) なので \((E,t)\) は AZ ample であり、
最後の主張はその生成条件から従う。
\end{proof}

\begin{example}[重み付き十二本矢と residual gerbe]
\label{ex:qa-weighted-not-h2}
\(S=\kk[x,y,z]\)、\(\deg x=\deg y=1\)、\(\deg z=2\) とし、
\(\mathcal X=\PP(1,1,2)\) を重み付き射影 \emph{stack} とする。
この環は global dimension 三、Gorenstein parameter 四の
noetherian AS 正則環である。
しかし
\[
 (\coh\mathcal X,L)\notin\HH\qquad(L\in\Pic(\mathcal X)).
\]
\end{example}
\begin{proof}
通常の三変数 Koszul 解消に各変数の重みを入れると
\[
 0\longrightarrow S(-4)
 \longrightarrow S(-3)^2\oplus S(-2)
 \longrightarrow S(-1)^2\oplus S(-2)
 \longrightarrow S\longrightarrow\kk\longrightarrow0
\]
を得る。これは AS 条件と Gorenstein parameter を与える。

\((\OO,\OO(1),\OO(2),\OO(3))\) の有限 End 代数を考える。
隣接する \(x_i,y_i\) が六本、距離二の \(z_0,z_1\) が二本であり、
最小関係は
\[
 \begin{aligned}
 x_1y_0-y_1x_0&=0,&x_2y_1-y_2x_1&=0,\\
 z_1x_0-x_2z_0&=0,&z_1y_0-y_2z_0&=0.
 \end{aligned}
\]
従って relation completion は十二本矢である。
前節の四写像表示では \(a,b\) はともに階数三で、核は交代テンソルの直線、
\(c,d\) は全射である。よってこの点は実際に
\(\mathfrak F_{\le12}\) の最小十二本矢 locus に属する。
この四対象が full strong であることも直接確認できる。
重み付き多項式環の局所コホモロジーから
\(H^1(\mathcal X,\OO(n))=0\) であり、
Serre 双対により \(n\ge-3\) では \(H^2(\mathcal X,\OO(n))=0\) である。
これで strong 性が従う。
三変数の sheaf Koszul 完全列は四つの連続する twist から
全ての twist を生成し、有限生成 graded \(S\)-加群の有限自由解消により
これらは \(D^b(\coh\mathcal X)\) 全体を生成する。

一方、重み付き Euler 列から
\(\omega_{\mathcal X}=\OO(-4)\) であり、\(t=-\otimes\OO(4)\) である。
点 \([0:0:1]\) の residual gerbe を
\(i:B\mu_2\hookrightarrow\mathcal X\) とし、
\(\chi\) を \(\mu_2\) の非自明な一次元表現とする。
任意の可逆層 \(L\) に対し、非零 coherent 対象
\(M_L=i_*(i^*L\otimes\chi)\) を取る。
\(\OO(4n)\) の gerbe 上の指標は自明なので、引戻しと直像の随伴により
\[
 \Hom(L(-4n),M_L)=
 H^0(B\mu_2,i^*\OO(4n)\otimes\chi)=0\qquad(n\ge0).
\]
特に標準 helix の四位相すべてで生成条件が破れる。
補題~\ref{lem:qa-antiserre-necessary} に反するので、主張を得る。
\end{proof}

\begin{example}[\(\Sigma_2\) 自身の構造対象]
\label{ex:qa-hirzebruch-not-h2}
\(Y=\Sigma_2\) の負断面を \(C\)、ファイバーを \(f\) とする。
\(C^2=-2\)、\(-K_Y=2C+4f\) なので、
\(\OO_Y(-K_Y)|_C\simeq\OO_C\) である。
任意の可逆層 \(L\) に対し、
\(M_L=i_*(i^*L\otimes\OO_C(-1))\ne0\) とおくと
\[
 \Hom\bigl(L\otimes\OO_Y(nK_Y),M_L\bigr)
   =H^0(C,\OO_C(-1))=0\qquad(n\ge0).
\]
同じ必要条件から \((\coh\Sigma_2,L)\notin\HH\) となる。
これは \emph{avatar の一点 Hilbert 曲面が \(\Sigma_2\) である}ことと、
元の Abel 圏が \(\coh\Sigma_2\) であることを混同してよいという意味ではない。
Hilbert 曲面の型だけから、対応する非可換圏の \(\HH\) 適合性は決まらない。
\end{example}

この二例は、十二本矢の全 locus、AS 正則性、あるいは
full strong foundation の存在だけで忘却の標的が \(\HHstack\) になる、
という主張を否定する。
適切な開部分が存在する可能性を否定するものではなく、
その開部分には構造対象の反 Serre 軌道に対する生成条件も必要であることを示す。

\subsection{標識、射、既存の表示スタック}

\(\HHstack\) への忘却では \(\OO\) を保持する。
従って、foundation の mutation が構造対象を別の対象へ移す場合には、
そのまま同じ pointed 対象の表示変更と扱えない。
QP の mutation が dg 導来同値を与えても、
指定した heart の完全同値と \(\OO\) の同型は別途必要である。
\cite[\texttt{cor:pointed-derived-exact}]{QuadricIntrinsic2026} は体上の \(\HH\) 対象の間の判定であり、
任意の相対 periodic family について同じ結論を自動的に与えるものではない。

同様に、消去可能な二本の矢 \(x,y\) と項 \(yx\) を付加した生の QP には、
\((x,y)\mapsto(\lambda x,\lambda^{-1}y)\) という自己同型がある。
消去後の表示ではこの自己同型は見えなくなる。
従って、生の安定化 QP groupoid の消去を、
射を補正しないままスタック同値とすることはできない。
有限 foundation の頂点スカラーの剛体化、
周期的環の次数スカラーの除去、
構造対象を保つ自己同値全体での同一視は、それぞれ異なる操作である。

現在の内在的研究との関係をまとめると、
\cite[\texttt{prop:foundation-stack}]{QuadricModuli2026} の有限二写像表示 \(\mathfrak F_0\) は、
一方の長さ二の積を可逆にした foundation chart である。

\begin{proposition}[幾何点における十本矢 chart の被覆]
\label{prop:qa-h2-normal-ten-chart}
体上の \(\HH\) の任意の対象は、
\(\mathfrak F_{\le10}^{(a)}\) に属する normal foundation を持つ。
これは幾何点での主張であり、任意の \(\HHstack\) の族における
foundation の相対的な存在・降下を結論するものではない。
\end{proposition}
\begin{proof}
\cite[\texttt{thm:normal-helix}]{QuadricIntrinsic2026} と
\cite[\texttt{prop:ruling-involution-mutation}]{QuadricModuli2026} により、
\((E_0,E_1,E_2,E_3)=(\OO,X,s\OO,Z)\) という foundation があり、
本節の \(b\) は同型、\(c\) は全射である。
\(E_0,E_1,E_2\) は rank-one torsion-free 対象である。
このような対象間の非零射は単射である。
実際、像が低い Gabriel--Krull 段階に入れば標的の torsion-free 性に反し、
従って rank-one 商では非零射、すなわち同型を誘導する。
その核は低い段階に入るが、始域の torsion-free 性によって零である。
従って二つの隣接する非零射の合成は零にならず、
\(\ker a\) は非零の分解可能テンソルを含まない。
\(b\) は同型なので他方の核は零であり、
\(c\) の全射性から定義~\ref{def:qa-foundation-open} の全条件を満たす。
\end{proof}

相対的な ruling 交換を \(\tau\) とすると、
\[
 \mathfrak Q_{\mathrm{fnd}}=[\mathfrak F_0/S_2]
\]
は有限型 Artin スタックであり、体上では
\(\HH(\kk)\) をその full subgroupoid として実現する
（\cite[\texttt{cor:h2-point-groupoid-foundation}]{QuadricModuli2026}）。
この幾何点の結果だけから、\(\mathfrak Q_{\mathrm{fnd}}\) 全体または
その未特定の開部分が \(\HHstack\) と同値であるとはいえない。

相対的には、同論文の反 Serre 認識によって全忠実な射
\[
 \HHstack\hookrightarrow
 \mathfrak P=[X_3/\PGL_9]
\]
がある。
ここで \(X_3\subset\operatorname{Gr}(56,81)\) は、
九生成元・五十六二次関係の環の三次成分が階数四十九となる
scheme-theoretic な局所閉 locus である。
\[
 X_{\mathrm{as}}=X_3\times_{\mathfrak P}\HHstack
\]
は不変 fppf 部分層として定義され、
\(\HHstack\simeq[X_{\mathrm{as}}/\PGL_9]\) が成立する。
\(X_{\mathrm{as}}\) の有限型・滑らか・既約な純次元83のスキームとしての表現可能性は
相対モジュライ論文の定理~\cite[\texttt{thm:global-h2-smooth-finite-type}]{QuadricModuli2026} と
系~\cite[\texttt{cor:global-xas-smooth-scheme}]{QuadricModuli2026} で証明されている。
ただし $X_{\mathrm{as}}\to X_3$ は一致根対で開ではなく、
局所閉性と明示的な偏極認識方程式は未解決である。
従って本稿の八・十・十二本矢の明示的な開部分は、
この相対認識問題への有限表示の候補を与えるが、
それ自体で内在的モジュライの代数性を解決するものではない。


\section{比較を完成させるための認識問題}
\label{sec:qa-remaining}

本稿の有限代数の開部分と、幾何的 avatar の開部分は、それぞれ独立に定義される。
両者を相対 AS 分類として同定するには、avatar から得る有限代数が
全周期にわたる相対 AS 代数を再構成し、逆向きの幾何構成と任意底変更に可換すること、
および同型の層と自己同型群が一致することを証明する必要がある。
体上の全単射はこの最後の同型層の比較を代替しない。

さらに $\HHstack$ へ進むには、選んだ構造対象に対する標準 heart、
相対 Serre 双対、豊富な exact 平方根、その同じ平方根の全軌道環の
正則正規 AS4 被覆を認識する必要がある。
これらを満たす部分には忘却射があるが、
その部分が十二本矢表示モジュライの開部分であることは別の命題である。
平方根や被覆を選んだ有限表示の incidence が存在することから、
それらの局所存在だけを要求する部分の開性や代数性を推論してはならない。

従って次の研究課題は、矢数だけを減らす追加の消去ではなく、
相対 AS 再構成と内在的認識の二つの比較を閉じることである。
どちらの場合も、marked な表示の自己同型から消える核、
構造対象の mutation による移動、非被約底での基底変更を同時に追跡する。
この順序により、有限型の表示空間を先に構成したことと、
内在的モジュライを表現したことの論理的な境界が明確になる。

\bibliographystyle{alpha}
\bibliography{references}
\end{document}
